\section{Actual cotangent, exterior and canonical line sheaves}

This Mathlib-only proof cone is ported from the actual sibling project
\texttt{BingWangEntropy}. Exact upstream source hashes, the original
source locations, the source commit and every namespace/import compatibility
diff are retained in the accompanying provenance manifest. The objects
are actual original scheme module sheaves. Their differential and
rank-one frame properties are proved from actual smoothness and
localization. No vanishing or uniform multiplier-generation theorem
is claimed by these constructions.

\begin{definition}[The literal relative exterior canonical sheaf]
\label{sc0-canonical-port:can:actual-relative-exterior-canonical-object}
\uses{sc0-canonical-port:ext:relative-exterior-cotangent-sheaf}
\lean{EnvelopeProperty.CanonicalGeometry.canonicalLineSheaf}
\leanok
For \(n\in\mathbb N\), let \(\mathcal K_{X/K,n}\)
be the actual sheafification of the presheaf
\(V\mapsto\bigwedge^n_{\Gamma(X,V)}\Omega_{\Gamma(X,V)/K}\),
with its actual semilinear differential restriction maps.
Thus its sheaf is literally the actual relative exterior
cotangent sheaf \(\Omega^{[n]}_{X/K}\). Smoothness below
produces its line atlas when the relative dimension is \(n\).
\end{definition}

\begin{theorem}[Basis-free exterior transport under an actual coefficient-ring isomorphism]
\label{sc0-canonical-port:can:basis-free-semilinear-exterior-transport}
\uses{sc0-canonical-port:ext:semilinear-exterior-map}
\lean{EnvelopeProperty.CanonicalGeometry.exteriorSemilinearMap_bijective_of_ringIso}
\leanok
A bijective semilinear map \(h:M\to N\) for an actual ring
isomorphism \(\rho:R\simeq S\) induces an actual exterior equivalence.
Its inverse is induced by the inverse semilinear map, without a basis.
\end{theorem}
\begin{proof}
Apply injectivity of \(h\) to test additivity of its inverse.
For \(s\in S\), the image under \(h\) of
\(\rho^{-1}(s)h^{-1}(y)\) is \(sy\), proving inverse semilinearity.
Apply the actual exterior construction to \(h\) and its inverse.
Their composites fix every pure wedge by the inverse identities.
Wedges span, so the composites are identities everywhere.
\end{proof}

\begin{theorem}[Raw relative exterior cotangent presheaves restrict along actual open immersions]
\label{sc0-canonical-port:can:raw-exterior-cotangent-open-restriction}
\uses{sc0-canonical-port:can:basis-free-semilinear-exterior-transport,sc0-canonical-port:geom:cotangent-presheaf-open-immersion-restriction,sc0-canonical-port:ext:relative-exterior-cotangent-sheaf}
\lean{EnvelopeProperty.CanonicalGeometry.cotangentOpenImmersionSemilinear,EnvelopeProperty.CanonicalGeometry.exteriorCotangentOpenImmersionComponent,EnvelopeProperty.CanonicalGeometry.exteriorCotangentOpenImmersionComponent_wedge,EnvelopeProperty.CanonicalGeometry.exteriorCotangentPresheafOpenImmersionMap,EnvelopeProperty.CanonicalGeometry.exteriorCotangentOpenImmersionComponent_bijective,EnvelopeProperty.CanonicalGeometry.exteriorCotangentPresheafOpenImmersionIso}
\leanok
For an actual open immersion \(g:U\to X\), the raw relative exterior
cotangent presheaf of \(f\circ g\) is naturally isomorphic to the
restricted raw exterior cotangent presheaf of \(f\), with actual
coefficient-ring scalar transport.
\end{theorem}
\begin{proof}
The actual cotangent-presheaf restriction isomorphism supplies
coefficient-ring isomorphisms and bijective semilinear differential
maps on every open. Apply basis-free exterior transport.
On pure wedges, both restriction composites are wedges of the equal
cotangent composites. Wedge spanning proves full naturality.
The pointwise inverse identities give the presheaf isomorphism.
\end{proof}

\begin{theorem}[Actual relative exterior cotangent sheaves restrict along open immersions]
\label{sc0-canonical-port:can:relative-exterior-cotangent-sheaf-open-restriction}
\uses{sc0-canonical-port:can:raw-exterior-cotangent-open-restriction,sc0-canonical-port:geom:sheafification-open-immersion-restriction}
\lean{EnvelopeProperty.CanonicalGeometry.exteriorCotangentSheafOpenImmersionIso}
\leanok
The actual relative exterior cotangent sheaf of \(f\) restricts
along \(g\) to that of \(f\circ g\).
\end{theorem}
\begin{proof}
The actual open-immersion sheafification comparison identifies
restricted exterior sheafification with sheafification of the
restricted raw exterior presheaf. Sheafify the preceding raw
presheaf isomorphism and compose the two actual comparisons.
\end{proof}

\begin{theorem}[A genuine standard-smooth chart identifies the global exterior sheaf with its actual coordinate exterior module]
\label{sc0-canonical-port:can:actual-smooth-chart-exterior-coordinate-comparison}
\uses{sc0-canonical-port:can:relative-exterior-cotangent-sheaf-open-restriction,sc0-canonical-port:geom:affine-chart-structural-constants,sc0-canonical-port:ext:actual-affine-exterior-sheaf-comparison}
\lean{EnvelopeProperty.CanonicalGeometry.standardSmoothAffineChartExteriorCotangentSheafIso}
\leanok
Let \(V\subset X\) be an actual affine open. Suppose its
actual structural ring map \(K\to A=\Gamma(X,V)\)
is standard smooth of relative dimension \(n\).
Restriction of \(\Omega^{[n]}_{X/K}\) along its native
spectrum chart is isomorphic to the associated sheaf of
\(\bigwedge^n_A\Omega_{A/K}\).
\end{theorem}
\begin{proof}
The genuine open-restriction comparison gives the actual
relative exterior sheaf of the inherited affine structure.
The actual structural-constant identity makes this the actual
coordinate structure of \(A\). Standard smoothness produces
an actual \(A\)-basis of \(\Omega_{A/K}\) with \(n\)
elements. The accepted finite-free affine exterior comparison,
with that produced basis, identifies the exterior sheaf with
the associated sheaf of the actual exterior module.
Composing these genuine isomorphisms gives the claimed map.
\end{proof}

\begin{theorem}[Genuine standard-smooth scheme charts give actual relative top-exterior frames]
\label{sc0-canonical-port:can:actual-smooth-chart-relative-top-frame}
\uses{sc0-canonical-port:can:relative-exterior-cotangent-sheaf-open-restriction,sc0-canonical-port:geom:smooth-scheme-standard-smooth-coordinate-cover,sc0-canonical-port:geom:affine-chart-structural-constants,sc0-canonical-port:ext:standard-smooth-actual-top-exterior-frame,sc0-canonical-port:can:actual-relative-exterior-canonical-object,sc0-canonical-port:can:actual-smooth-chart-exterior-coordinate-comparison}
\lean{EnvelopeProperty.CanonicalGeometry.standardSmoothCanonicalAffineChartFrame,EnvelopeProperty.CanonicalGeometry.smooth_canonicalLineSheaf_local_frame}
\leanok
For an actual smooth relative-dimension-\(n\) structure
\(f:X\to\operatorname{Spec}K\), each genuine standard-smooth affine
chart supplies an actual local frame of \(\Omega_{X/K}^{[n]}\).
\end{theorem}
\begin{proof}
Restrict the actual relative exterior sheaf through its genuine
open-immersion comparison. The actual structural-constant identity
identifies the inherited affine structure with its coordinate
\(K\)-algebra structure. The latter is actually standard smooth of
relative dimension \(n\), so its proved relative top-exterior frame
applies, with a differential basis produced by smoothness.
Transport that frame through the genuine affine chart isomorphism
and the actual open-restriction comparison. This gives the actual
local frame on the original open.
\end{proof}

\begin{theorem}[Actual smooth top-exterior frames construct a genuine global line atlas]
\label{sc0-canonical-port:can:actual-smooth-relative-exterior-line-atlas}
\uses{sc0-canonical-port:can:actual-smooth-chart-relative-top-frame,sc0-canonical-port:geom:smooth-scheme-standard-smooth-coordinate-cover}
\lean{EnvelopeProperty.CanonicalGeometry.smoothCanonicalLineAtlas,EnvelopeProperty.CanonicalGeometry.smoothCanonicalLineBundle,EnvelopeProperty.CanonicalGeometry.smoothCanonicalLineBundle_sheaf}
\leanok
The actual relative top exterior cotangent sheaf of a smooth
relative-dimension-\(n\) scheme over \(K\) has a produced genuine
line atlas. Its intrinsic line object has sheaf literally
\(\Omega_{X/K}^{[n]}\).
\end{theorem}
\begin{proof}
For each actual point \(x\), choose the produced standard-smooth
affine open \(V_x\) containing it and its produced actual frame.
These genuine opens cover, since \(x\in V_x\) for every \(x\).
The point-indexed cover and its actual frames form a line atlas
on the actual relative exterior sheaf. Packaging that actual sheaf
and atlas gives the intrinsic line object. No frame or line-atlas
conclusion was assumed as an input.
\end{proof}

\begin{definition}[Actual affine scalar actions and coordinate constants]
\label{sc0-canonical-port:geom:affine-coordinate-scalar-actions}
\lean{EnvelopeProperty.CanonicalGeometry.affineCoordinateMap,EnvelopeProperty.CanonicalGeometry.affineCoordinateMap_naturality,EnvelopeProperty.CanonicalGeometry.specPresheafToAModules,EnvelopeProperty.CanonicalGeometry.specPresheafAModuleMap,EnvelopeProperty.CanonicalGeometry.specPresheafToAModules_sheaf}
\leanok
For a commutative ring \(A\) and \(X=\operatorname{Spec}A\),
the actual map \(A\to\mathcal O_X(U)\) is the canonical
global-coordinate identification followed by restriction.
These maps commute with restriction. Every actual structure-module
presheaf can consequently be regarded as a presheaf of actual
\(A\)-modules by restricting scalars through these maps;
its morphisms give actual \(A\)-linear presheaf maps.
For an actual module sheaf, this action is the usual affine
coordinate action.
\end{definition}

\begin{theorem}[Natural affine coordinate-linear maps into a sheaf are regular-function linear]
\label{sc0-canonical-port:geom:affine-coordinate-linear-map-upgrade}
\uses{sc0-canonical-port:geom:affine-coordinate-scalar-actions}
\lean{EnvelopeProperty.CanonicalGeometry.affinePresheafHomOfAModuleHom}
\leanok
If \(P\) is an actual structure-module presheaf and
\(N\) an actual structure-module sheaf on
\(\operatorname{Spec}A\), every natural affine
\(A\)-linear map \(P\to N\) is linear over each actual
regular-function ring \(\mathcal O_X(U)\).
\end{theorem}
\begin{proof}
On a principal open \(D(r)\), both section modules have
their actual \(A_r\)-actions. Any \(A\)-linear map
between them is \(A_r\)-linear: it commutes with
multiplication by \(r\), and hence with its inverse,
so with every fraction \(a/r^k\).
For a regular function \(t\) on an arbitrary \(U\),
cover \(U\) by principal opens contained in it.
Naturality restricts the claimed identity
\(h_U(ts)=t h_U(s)\) to the just-proved identity on
each such open. Target sheaf separation makes the
identity hold on \(U\), proving the upgrade.
\end{proof}

\begin{theorem}[Actual maps on the principal-open basis extend uniquely]
\label{sc0-canonical-port:geom:affine-principal-basis-extension}
\uses{sc0-canonical-port:geom:affine-coordinate-linear-map-upgrade}
\lean{EnvelopeProperty.CanonicalGeometry.affineBasisExtend,EnvelopeProperty.CanonicalGeometry.affineBasisExtend_apply,EnvelopeProperty.CanonicalGeometry.affineBasisHomExt}
\leanok
An actual natural \(A\)-linear map on the principal-open
basis from an actual structure-module presheaf \(P\)
to an actual module sheaf \(N\) extends uniquely to an
actual structure-module presheaf map \(P\to N\).
The extension equals its prescribed map on every
principal open. Two such presheaf maps are equal if
their principal-open components are equal.
\end{theorem}
\begin{proof}
For \(s\in P(U)\), restrict to each principal open
\(D(r)\subset U\) and apply its prescribed map.
On intersections \(D(r)\cap D(t)=D(rt)\), naturality
makes their further restrictions equal. The actual
sheaf axiom therefore glues them to a unique section
of \(N(U)\). This construction commutes with
restriction, addition, and the affine \(A\)-action,
since each identity holds locally and gluing is unique.
The preceding theorem makes it regular-function linear.
On a principal open the prescribed image already has
all these restrictions, so uniqueness gives the stated
component equality. The same separation argument on
this basis proves uniqueness and the final equality
criterion.
\end{proof}

\begin{theorem}[Actual localized maps give a natural principal-open comparison]
\label{sc0-canonical-port:geom:affine-localized-basis-comparison}
\uses{sc0-canonical-port:geom:affine-coordinate-scalar-actions}
\lean{EnvelopeProperty.CanonicalGeometry.affineLocalizedBasisHom,EnvelopeProperty.CanonicalGeometry.affineLocalizedBasisIso}
\leanok
Let \(M\) be an actual \(A\)-module and \(P\) an
actual structure-module presheaf on \(\operatorname{Spec}A\).
Suppose actual \(A\)-linear maps \(\phi_U:M\to P(U)\)
commute with restrictions and, for every \(r\in A\),
\(\phi_{D(r)}\) makes \(P(D(r))\) the actual module
localization of \(M\) at the powers of \(r\).
Then there is a natural isomorphism on the actual
principal-open basis
\(P(D(r))\simeq\widetilde M(D(r))\), taking
\(\phi_{D(r)}(m)\) to the actual localized image of \(m\).
\end{theorem}
\begin{proof}
Both modules on a principal open are actual localizations
of \(M\) at the powers of \(r\). The localization
universal property gives unique inverse maps fixing
the images of every \(m\). For an inclusion of
principal opens, each restriction composite fixes the
same original images, by the naturality of \(\phi\)
and the actual associated-sheaf restriction. Uniqueness
of localized maps makes these composites equal. This
proves naturality, as well as the asserted generator
formula and actual basis isomorphism.
\end{proof}

\begin{theorem}[Actual affine sheafification comparison maps]
\label{sc0-canonical-port:geom:affine-localized-sheafification-maps}
\uses{sc0-canonical-port:geom:affine-localized-basis-comparison,sc0-canonical-port:geom:affine-principal-basis-extension}
\lean{EnvelopeProperty.CanonicalGeometry.affinePresheafSheafification,EnvelopeProperty.CanonicalGeometry.affinePresheafSheafificationUnit,EnvelopeProperty.CanonicalGeometry.affineLocalizedPresheafHom,EnvelopeProperty.CanonicalGeometry.affineLocalizedSheafificationHom,EnvelopeProperty.CanonicalGeometry.affineLocalizedSheafificationInv,EnvelopeProperty.CanonicalGeometry.affineLocalizedSheafificationHom_fac,EnvelopeProperty.CanonicalGeometry.affineLocalizedSheafificationHom_unit_apply,EnvelopeProperty.CanonicalGeometry.affineLocalizedSheafificationInv_apply}
\leanok
Under the actual localized-map hypotheses of
Theorem~\ref{sc0-canonical-port:geom:affine-localized-basis-comparison}, let
\(q:P\to P^\#\) be the actual module sheafification
unit, and \(\beta_r:P(D(r))\simeq\widetilde M(D(r))\)
the proved comparison. There are actual sheaf maps
\(h:P^\#\to\widetilde M\) and
\(j:\widetilde M\to P^\#\) satisfying on each
principal open
\[
 h(q(x))=\beta_r(x),\qquad
 j(y)=q(\beta_r^{-1}(y)).
\]
The first factors the uniquely extended presheaf comparison.
\end{theorem}
\begin{proof}
Extend the actual basis comparison to a presheaf map
\(P\to\widetilde M\) by
Theorem~\ref{sc0-canonical-port:geom:affine-principal-basis-extension}.
Sheafification's unique extension property gives
\(h\) factoring this map through \(q\), proving
its displayed formula on the basis.
Compose the inverse basis comparison with the actual
unit \(q\). These maps are natural on the basis,
so the same extension theorem produces a map
\(\widetilde M\to P^\#\). Both objects are actual
sheaves; this is the actual sheaf map \(j\), and
the extension's basis formula is exactly the second
displayed identity.
\end{proof}

\begin{theorem}[Actual localized maps produce the global affine sheaf isomorphism]
\label{sc0-canonical-port:geom:affine-localized-sheafification-isomorphism}
\uses{sc0-canonical-port:geom:affine-localized-sheafification-maps,sc0-canonical-port:geom:affine-principal-basis-extension}
\lean{EnvelopeProperty.CanonicalGeometry.affineLocalizedSheafificationHom_inv_id,EnvelopeProperty.CanonicalGeometry.affineLocalizedSheafificationInv_hom_id,EnvelopeProperty.CanonicalGeometry.affineLocalizedSheafificationIso}
\leanok
Under those same explicit localized-map hypotheses,
the actual maps \(h,j\) are inverse, giving an actual
\(\mathcal O_X\)-module sheaf isomorphism
\(P^\#\simeq\widetilde M\).
\end{theorem}
\begin{proof}
On every principal open, the two formulas of the
preceding theorem give \(h(j(y))=y\). Equality on the
basis gives \(h\circ j=\mathrm{id}\) on the actual
associated sheaf.
They also give \(j(h(q(x)))=q(x)\) on each principal
open. The basis extension uniqueness makes the
presheaf composites on \(P\) equal everywhere.
The sheafification universal property determines
maps out of \(P^\#\) by their composite with \(q\),
so \(j\circ h=\mathrm{id}\) on \(P^\#\).
Thus the two maps are actual inverses and give the
asserted sheaf isomorphism.
\end{proof}

\begin{definition}[The actual affine relative cotangent objects]
\label{sc0-canonical-port:geom:actual-affine-cotangent-objects}
\uses{sc0-canonical-port:geom:cotangent-presheaf,sc0-canonical-port:geom:affine-coordinate-scalar-actions}
\lean{EnvelopeProperty.CanonicalGeometry.affineBaseRingMap,EnvelopeProperty.CanonicalGeometry.affineStructuralMap,EnvelopeProperty.CanonicalGeometry.affineCotangentModule,EnvelopeProperty.CanonicalGeometry.affineCotangentTilde,EnvelopeProperty.CanonicalGeometry.affineCotangentAPresheaf,EnvelopeProperty.CanonicalGeometry.affineBaseRingMap_coordinateMap}
\leanok
For a field \(K\) and a commutative \(K\)-algebra \(A\),
the actual structural morphism is
\(\operatorname{Spec}A\to\operatorname{Spec}K\),
the coordinate relative differential module is
\(M=\Omega_{A/K}\), and its actual associated sheaf
is \(\widetilde M\). The actual cotangent presheaf
can be regarded as an affine \(A\)-module presheaf.
Its structural constants on each open are the actual
composite \(K\to A\to\mathcal O_X(U)\).
\end{definition}

\begin{theorem}[Actual affine differential maps satisfy the localized-module hypotheses]
\label{sc0-canonical-port:geom:actual-affine-differential-localized-maps}
\uses{sc0-canonical-port:geom:actual-affine-cotangent-objects,sc0-canonical-port:geom:affine-cotangent-localization}
\lean{EnvelopeProperty.CanonicalGeometry.affineCotangentToOpen,EnvelopeProperty.CanonicalGeometry.affineCotangentToOpen_d,EnvelopeProperty.CanonicalGeometry.affineCotangentToOpen_naturality,EnvelopeProperty.CanonicalGeometry.affineCotangentToOpen_isLocalized}
\leanok
The actual maps
\(\phi_U:\Omega_{A/K}\to\Omega_{\mathcal O_X(U)/K}\)
send \(da\) to the actual differential of its restricted
coordinate function. They commute with restriction.
On every actual principal open they make the target
the actual module localization of \(\Omega_{A/K}\).
\end{theorem}
\begin{proof}
The commutative actual structural-constant square induces
these maps by the Kähler universal property, with the
stated formula on universal differential generators.
For restriction to a smaller open, the two candidate
composites agree on every \(da\), because the actual
coordinate functions and actual derivation commute
with restriction. The generators span, so the composites
agree. On \(D(r)\), the target is the actual
differential module of \(A_r\); its canonical
localization description is
Theorem~\ref{sc0-canonical-port:geom:affine-cotangent-localization}.
This proves the actual localized-module property
of the displayed canonical map.
\end{proof}

\begin{theorem}[The actual affine cotangent sheaf is the associated coordinate differential sheaf]
\label{sc0-canonical-port:geom:actual-affine-cotangent-sheaf-comparison}
\uses{sc0-canonical-port:geom:actual-affine-differential-localized-maps,sc0-canonical-port:geom:affine-localized-sheafification-isomorphism,sc0-canonical-port:geom:cotangent-sheaf}
\lean{EnvelopeProperty.CanonicalGeometry.affineCotangentSheafIso,EnvelopeProperty.CanonicalGeometry.affineCotangentModuleIso}
\leanok
There is an actual sheaf isomorphism on
\(X=\operatorname{Spec}A\),
\[
 \Omega_{X/K}\simeq\widetilde{\Omega_{A/K}}.
\]
The coordinate differential object is the actual
Kähler module for the original \(K\)-algebra structure.
\end{theorem}
\begin{proof}
The actual localized-map hypotheses are produced by
Theorem~\ref{sc0-canonical-port:geom:actual-affine-differential-localized-maps}.
Apply Theorem~\ref{sc0-canonical-port:geom:affine-localized-sheafification-isomorphism}
to the actual cotangent presheaf and actual coordinate
differential module. Its sheafification is exactly
\(\Omega_{X/K}\) by definition. The categorical
coordinate differential construction uses the original
ring map \(K\to A\); the induced algebra action is
that same action. Thus its generators and relations
are exactly those of the original \(\Omega_{A/K}\),
giving the asserted identification and actual sheaf
comparison.
\end{proof}

\begin{theorem}[An actual standard-smooth presentation gives a free cotangent sheaf]
\label{sc0-canonical-port:geom:standard-smooth-actual-cotangent-sheaf-frame}
\uses{sc0-canonical-port:geom:standard-smooth-differential-basis,sc0-canonical-port:geom:actual-affine-cotangent-sheaf-comparison}
\lean{EnvelopeProperty.CanonicalGeometry.standardSmoothAffineCotangentModuleFrame,EnvelopeProperty.CanonicalGeometry.standardSmoothAffineCotangentSheafFrame,EnvelopeProperty.CanonicalGeometry.standardSmoothAffineCotangentSheaf_locallyFree}
\leanok
If \(A\) is actually standard smooth of relative
dimension \(n\) over \(K\), its actual cotangent
sheaf on \(\operatorname{Spec}A\) has a sheaf frame
to the free structure-module sheaf on \(n\) generators.
In particular this actual cotangent sheaf is locally free.
\end{theorem}
\begin{proof}
The actual standard-smooth presentation produces a
coordinate differential basis of \(n\) vectors by
Theorem~\ref{sc0-canonical-port:geom:standard-smooth-differential-basis}.
Its coordinate map is an actual module frame to \(A^n\).
Apply the actual associated-sheaf functor to that
linear isomorphism; localizing its basis maps on
principal opens gives inverse sheaf maps to the free
structure-module sheaf. Compose with the actual
cotangent comparison of the preceding theorem.
This gives the stated actual free sheaf frame, which
itself supplies the local bases required for native
local freeness.
\end{proof}

\begin{definition}[Structural constants on every open set]
\label{sc0-canonical-port:geom:structural-constant-presheaf}
\lean{EnvelopeProperty.CanonicalGeometry.constantRingPresheaf,EnvelopeProperty.CanonicalGeometry.schemeConstantPresheafMap}
\leanok
Let \(K\) be a field, let \(X\) be a scheme, and let
\(f:X\to\operatorname{Spec}K\) be its structural morphism.
Write \(\underline K\) for the constant presheaf of rings.
The structural constant map on an open set \(U\subset X\) is
\[
 K\simeq\Gamma(\operatorname{Spec}K,\mathcal O)
 \xrightarrow{f^\#}\Gamma(X,\mathcal O_X)
 \longrightarrow\mathcal O_X(U).
\]
These maps define a morphism \(\underline K\to\mathcal O_X\)
of presheaves of rings.
\end{definition}

\begin{definition}[The objectwise relative differential presheaf]
\label{sc0-canonical-port:geom:cotangent-presheaf}
\uses{sc0-canonical-port:geom:structural-constant-presheaf}
\lean{EnvelopeProperty.CanonicalGeometry.cotangentPresheaf,EnvelopeProperty.CanonicalGeometry.cotangentPresheafDerivation}
\leanok
For each open \(U\), put
\(P(U)=\Omega_{\mathcal O_X(U)/K}\), with its actual
\(\mathcal O_X(U)\)-module structure. If \(V\subset U\), the
restriction is the map induced by \(\mathcal O_X(U)\to\mathcal
O_X(V)\); on differentials it sends \(da\) to
\(d(a|_V)\). The maps \(a\mapsto da\) form a derivation
\(d_P:\mathcal O_X\to P\) over \(\underline K\).
\end{definition}

\begin{theorem}[The universal derivation of the differential presheaf]
\label{sc0-canonical-port:geom:cotangent-presheaf-universal}
\uses{sc0-canonical-port:geom:cotangent-presheaf}
\lean{EnvelopeProperty.CanonicalGeometry.cotangentPresheafUniversal,EnvelopeProperty.CanonicalGeometry.cotangentPresheafHomEquiv}
\leanok
For every presheaf \(N\) of \(\mathcal O_X\)-modules,
postcomposition with \(d_P\) is a bijection between
\(\mathcal O_X\)-linear presheaf maps \(P\to N\) and
\(\underline K\)-derivations \(\mathcal O_X\to N\).
\end{theorem}
\begin{proof}
Given a derivation \(D\), the universal property of
\(\Omega_{\mathcal O_X(U)/K}\) gives a unique
\(\mathcal O_X(U)\)-linear map \(h_U\) satisfying
\(h_U(da)=D_U(a)\). For \(V\subset U\), the two candidate
restriction composites agree on every generator \(da\), because
\(D\) commutes with restriction. They therefore agree on the whole
differential module. Thus \((h_U)\) is a presheaf map.
Its composite with \(d_P\) is \(D\). Uniqueness on every open
set gives uniqueness as a presheaf map. Conversely, composing any
linear presheaf map with \(d_P\) is a derivation, so these
constructions are inverse.
\end{proof}

\begin{definition}[The actual relative cotangent sheaf]
\label{sc0-canonical-port:geom:cotangent-sheaf}
\uses{sc0-canonical-port:geom:cotangent-presheaf}
\lean{EnvelopeProperty.CanonicalGeometry.cotangentSheaf,EnvelopeProperty.CanonicalGeometry.cotangentSheafificationMap}
\leanok
The relative cotangent sheaf \(\Omega_{X/K}\) is the
sheafification of \(P\) as a presheaf of
\(\mathcal O_X\)-modules. Write
\(\eta:P\to\Omega_{X/K}\) for its canonical map.
\end{definition}

\begin{definition}[Differentiation into the cotangent sheaf]
\label{sc0-canonical-port:geom:cotangent-sheaf-derivation}
\uses{sc0-canonical-port:geom:cotangent-presheaf,sc0-canonical-port:geom:cotangent-sheaf}
\lean{EnvelopeProperty.CanonicalGeometry.cotangentDerivation}
\leanok
Define \(d_X=\eta\circ d_P\). This is a derivation of actual
regular functions over the structural constant presheaf, with
values in the underlying presheaf of \(\Omega_{X/K}\).
\end{definition}

\begin{lemma}[Differentiation laws and restriction]
\label{sc0-canonical-port:geom:cotangent-derivation-laws}
\uses{sc0-canonical-port:geom:cotangent-sheaf-derivation}
\lean{EnvelopeProperty.CanonicalGeometry.cotangentDerivation_add,EnvelopeProperty.CanonicalGeometry.cotangentDerivation_mul,EnvelopeProperty.CanonicalGeometry.cotangentDerivation_constant,EnvelopeProperty.CanonicalGeometry.cotangentDerivation_restrict}
\leanok
For regular functions \(a,b\) on \(U\) and a structural
constant \(c\in K\),
\[
 d_X(a+b)=d_Xa+d_Xb,\qquad
 d_X(ab)=a\,d_Xb+b\,d_Xa,\qquad d_X(c)=0.
\]
For \(V\subset U\), one also has
\(d_X(a|_V)=(d_Xa)|_V\).
\end{lemma}
\begin{proof}
The same identities hold for the universal derivation \(d_P\)
on each differential module. Apply the linear map \(\eta_U\)
to the additive and product identities, and apply it to the zero
identity for constants. Naturality of both \(d_P\) and
\(\eta\) gives the restriction identity.
\end{proof}

\begin{theorem}[The actual cotangent sheaf represents derivations]
\label{sc0-canonical-port:geom:cotangent-sheaf-universal}
\uses{sc0-canonical-port:geom:cotangent-presheaf-universal,sc0-canonical-port:geom:cotangent-sheaf,sc0-canonical-port:geom:cotangent-sheaf-derivation}
\lean{EnvelopeProperty.CanonicalGeometry.cotangentSheafHomEquiv,EnvelopeProperty.CanonicalGeometry.cotangentSheafHomEquiv_apply,EnvelopeProperty.CanonicalGeometry.cotangentSheaf_derivation_factorization}
\leanok
For every sheaf \(N\) of \(\mathcal O_X\)-modules and every
relative derivation \(D:\mathcal O_X\to N\), there is a unique
sheaf morphism \(h:\Omega_{X/K}\to N\) with
\(D=h\circ d_X\). Consequently, postcomposition with \(d_X\)
is a bijection from \(\operatorname{Hom}_{\mathcal O_X}
(\Omega_{X/K},N)\) to the relative derivations into \(N\).
\end{theorem}
\begin{proof}
Theorem~\ref{sc0-canonical-port:geom:cotangent-presheaf-universal} gives a unique
presheaf map \(P\to N\) inducing \(D\). The universal
property of sheafification gives a unique sheaf map \(h\)
whose composite with \(\eta\) is this map. It satisfies
\(h\circ d_X=D\). If another sheaf map has that property,
its composite with \(\eta\) agrees with the unique presheaf
map on the generators \(da\), hence everywhere. Uniqueness
in sheafification then makes the sheaf maps equal. This also
identifies the stated bijection explicitly as postcomposition.
\end{proof}

\begin{lemma}[The actual module on an open set]
\label{sc0-canonical-port:geom:cotangent-open-module}
\uses{sc0-canonical-port:geom:structural-constant-presheaf,sc0-canonical-port:geom:cotangent-presheaf}
\lean{EnvelopeProperty.CanonicalGeometry.cotangentOpenAlgebra,EnvelopeProperty.CanonicalGeometry.cotangentPresheafSectionsModule,EnvelopeProperty.CanonicalGeometry.cotangentPresheafSectionsEquiv}
\leanok
Give \(A=\mathcal O_X(U)\) the \(K\)-algebra structure
defined by the structural constants. Then the presheaf section
module \(P(U)\) is canonically \(A\)-linearly isomorphic to
the actual module \(\Omega_{A/K}\).
\end{lemma}
\begin{proof}
The construction of \(P\) takes the relative differential module
of the structural constant ring map on each open set. Its ring
action is multiplication by actual elements of \(A\). Thus the
two modules have the same generators \(da\), the same additivity,
Leibniz, and constant relations, and the same \(A\)-action.
The identity on generators defines inverse \(A\)-linear maps.
\end{proof}

\begin{theorem}[Standard smoothness gives objectwise differential freeness]
\label{sc0-canonical-port:geom:cotangent-presheaf-smooth-rank}
\uses{sc0-canonical-port:geom:cotangent-open-module,sc0-canonical-port:geom:standard-smooth-differential-basis}
\lean{EnvelopeProperty.CanonicalGeometry.cotangentPresheaf_free_of_standardSmooth,EnvelopeProperty.CanonicalGeometry.cotangentPresheaf_rank_of_standardSmooth}
\leanok
If the actual structural algebra \(A=\mathcal O_X(U)\) is
standard smooth over \(K\), then \(P(U)\) is free over \(A\).
If \(A\ne0\) and this algebra is standard smooth of relative
dimension \(n\), its module rank is \(n\).
These are statements about the presheaf module on the specified
open set.
\end{theorem}
\begin{proof}
The standard-smooth presentation produces the basis of
Lemma~\ref{sc0-canonical-port:geom:standard-smooth-differential-basis}. For standard
smoothness without a specified dimension, its finite submersive
presentation gives a basis indexed by the unselected coordinates.
Transport that basis through Lemma~\ref{sc0-canonical-port:geom:cotangent-open-module}
to prove freeness. In relative dimension \(n\), the produced
basis has \(n\) elements. Since \(A\ne0\), the rank of a
finite free module is the cardinality of its basis, giving \(n\).
\end{proof}

\begin{lemma}[Structural constants on an actual affine scheme]
\label{sc0-canonical-port:geom:affine-structural-constants}
\uses{sc0-canonical-port:geom:structural-constant-presheaf}
\lean{EnvelopeProperty.CanonicalGeometry.affine_schemeConstantPresheafMap,EnvelopeProperty.CanonicalGeometry.affineCotangentPresheafSectionsEquiv}
\leanok
Let \(A\) be a commutative \(K\)-algebra and
\(X=\operatorname{Spec}A\) with its induced structural morphism.
On every open \(U\subset X\), the structural constant map is
the actual composite \(K\to A\to\mathcal O_X(U)\).
With this algebra structure, \(P(U)\) is canonically
\(\mathcal O_X(U)\)-linearly isomorphic to
\(\Omega_{\mathcal O_X(U)/K}\).
\end{lemma}
\begin{proof}
On the full affine open, the canonical identification
\(A\simeq\Gamma(\operatorname{Spec}A,\mathcal O)\) is natural
in ring maps. Thus the map induced by
\(\operatorname{Spec}A\to\operatorname{Spec}K\) is exactly
\(K\to A\). Restricting to \(U\) gives the displayed
composite. Now apply Lemma~\ref{sc0-canonical-port:geom:cotangent-open-module} with
this identified structural algebra.
\end{proof}

\begin{theorem}[Principal-open differential localization]
\label{sc0-canonical-port:geom:affine-cotangent-localization}
\uses{sc0-canonical-port:geom:affine-structural-constants,sc0-canonical-port:geom:etale-differential-basis}
\lean{EnvelopeProperty.CanonicalGeometry.affineCotangentLocalizationEquiv,EnvelopeProperty.CanonicalGeometry.affineCotangentPresheafLocalizationEquiv,EnvelopeProperty.CanonicalGeometry.affineCotangentLocalizationEquiv_tmul_D}
\leanok
For \(r\in A\), put
\(B=\mathcal O_{\operatorname{Spec}A}(D(r))\), with its actual
localization map \(A\to B\). There are canonical
\(B\)-linear isomorphisms
\[
 B\otimes_A\Omega_{A/K}\simeq\Omega_{B/K}\simeq P(D(r)).
\]
The first sends \(b\otimes da\) to \(b\,d(\overline a)\),
where \(\overline a\) is the image of \(a\) in \(B\).
\end{theorem}
\begin{proof}
The actual ring of regular functions on \(D(r)\) is the
localization of \(A\) at the powers of \(r\). For any
\(B\)-module, a \(K\)-derivation of \(A\) extends uniquely
to this ring by the formula
\(D(a/r^m)=r^{-m}D(a)-m a r^{-m-1}D(r)\).
This formula respects localization equality: multiplying a relation
by a sufficiently high power of \(r\), differentiating it, and
then inverting \(r\) makes the corresponding derivative relation
zero. The Leibniz rule is immediate after expressing both terms
with a common denominator. Uniqueness follows by differentiating
\(rr^{-1}=1\).
Consequently the universal modules for these derivations are
\(B\otimes_A\Omega_{A/K}\) and \(\Omega_{B/K}\), yielding
inverse maps. Their defining rule on universal differentials gives
the asserted pure-tensor formula. The second isomorphism is
Lemma~\ref{sc0-canonical-port:geom:affine-structural-constants} on this open set.
\end{proof}

\begin{definition}[Actual open-immersion module-presheaf restriction]
\label{sc0-canonical-port:geom:open-immersion-module-presheaf-restriction}
\lean{EnvelopeProperty.CanonicalGeometry.openImmersionStructurePresheafMap,EnvelopeProperty.CanonicalGeometry.openImmersionPresheafRestrict}
\leanok
For an actual open immersion \(g:Y\to X\), the
actual coefficient map on an open \(W\subset Y\)
is the inverse of the native isomorphism
\(\Gamma(X,gW)\simeq\Gamma(Y,W)\).
Restrict a structure-module presheaf \(P\) using its
actual values on \(gW\), with the \(\Gamma(Y,W)\)-action
transported through that inverse. Restrictions are the
actual restrictions of \(P\).
\end{definition}

\begin{lemma}[An actual coefficient-ring isomorphism induces a bijective differential map]
\label{sc0-canonical-port:geom:kaehler-coefficient-isomorphism-bijective}
\uses{sc0-canonical-port:geom:cotangent-open-module}
\lean{EnvelopeProperty.CanonicalGeometry.kaehlerCoefficientMap_bijective}
\leanok
Let \(R,B,C\) be commutative rings, let \(R\to B\)
and \(R\to C\) be actual structure maps, and let
\(e:B\simeq C\) be an actual ring isomorphism commuting
with those structure maps. The canonical semilinear map
\(\Omega_{B/R}\to\Omega_{C/R}\) induced by \(e\)
is bijective.
\end{lemma}
\begin{proof}
The universal derivation induces the map \(db\mapsto d(e(b))\).
The inverse ring isomorphism commutes with the same constants,
so it induces \(dc\mapsto d(e^{-1}(c))\). Each composite
fixes all universal differentials, which generate the
corresponding module over its coefficient ring. Semilinearity
therefore makes each composite the identity, proving bijectivity.
\end{proof}

\begin{lemma}[Actual open-immersion coefficient isomorphisms preserve structural constants]
\label{sc0-canonical-port:geom:open-immersion-constant-squares}
\uses{sc0-canonical-port:geom:open-immersion-module-presheaf-restriction}
\lean{EnvelopeProperty.CanonicalGeometry.openImmersion_constantSquare,EnvelopeProperty.CanonicalGeometry.openImmersion_constantSquare_inv}
\leanok
For an actual \(f:X\to\operatorname{Spec}K\), an actual
open immersion \(g:Y\to X\), and an open \(W\subset Y\),
the actual isomorphism \(\Gamma(X,gW)\simeq\Gamma(Y,W)\)
sends the structural constants for \(f\) to those for
\(f\circ g\). Its inverse sends them back.
\end{lemma}
\begin{proof}
Structural constants are obtained from the actual field
spectrum global-section isomorphism, followed by the structure
morphism's global-section map and then restriction to the
specified open. Naturality of the structure-sheaf map of \(g\)
identifies the restricted composite with that for \(f\circ g\).
This proves the forward equality for every constant. Applying
the inverse coefficient isomorphism and its inverse identity
gives the backward equality.
\end{proof}

\begin{lemma}[The actual restricted sheafification unit remains locally bijective]
\label{sc0-canonical-port:geom:open-immersion-restricted-sheafification-unit}
\uses{sc0-canonical-port:geom:open-immersion-module-presheaf-restriction,sc0-canonical-port:geom:sheafification-open-restriction}
\lean{EnvelopeProperty.CanonicalGeometry.openImmersionSheafificationMap,EnvelopeProperty.CanonicalGeometry.openImmersionSheafificationMap_underlying,EnvelopeProperty.CanonicalGeometry.openImmersionSheafificationMap_locally_bijective}
\leanok
For any actual structure-module presheaf \(P\) on \(X\)
and actual open immersion \(g:Y\to X\), the restricted
actual unit \(P\to P^{\mathrm{sh}}\) is a natural map
from the restricted presheaf to the restriction of its
sheafification. Its underlying additive-presheaf map is
literally the restricted unit. It is locally injective
and locally surjective on \(Y\).
\end{lemma}
\begin{proof}
On each open use the actual unit component on its image open.
Transport of the scalar action through the inverse coefficient
isomorphism preserves its linearity, and the unit's naturality
preserves restriction. This proves the asserted underlying map.
The actual unit is locally bijective on \(X\). Pull any
cover giving local lifts, or giving equality of two unit
preimages, back through the open immersion. These pulled-back
opens cover the relevant open of \(Y\), and the same actual
component identities hold there. Hence both local surjectivity
and local injectivity hold on \(Y\).
\end{proof}

\begin{theorem}[Actual cotangent presheaves restrict along open immersions]
\label{sc0-canonical-port:geom:cotangent-presheaf-open-immersion-restriction}
\uses{sc0-canonical-port:geom:cotangent-sheaf,sc0-canonical-port:geom:cotangent-open-module,sc0-canonical-port:geom:kaehler-coefficient-isomorphism-bijective,sc0-canonical-port:geom:open-immersion-constant-squares,sc0-canonical-port:geom:open-immersion-module-presheaf-restriction}
\lean{EnvelopeProperty.CanonicalGeometry.cotangentPresheafOpenImmersionMap,EnvelopeProperty.CanonicalGeometry.cotangentPresheafOpenImmersionIso}
\leanok
For an actual open immersion \(g:U\to X\), restriction
of the actual cotangent presheaf of \(f\), with scalars
transported through the actual open-immersion coefficient-ring
isomorphisms, is naturally isomorphic to the actual cotangent
presheaf of \(f\circ g\).
\end{theorem}
\begin{proof}
On each open of \(U\), the actual open immersion identifies
its regular-function ring with that on its image open in \(X\).
Structural constants commute with these isomorphisms because
the structure morphism of \(U\) is \(f\circ g\).
The universal property of K\"ahler differentials therefore gives
the semilinear differential map induced by each coefficient-ring
isomorphism. The differential map induced by the inverse ring
isomorphism is its inverse: both composites send every universal
differential \(da\) to \(da\), and those differentials generate
the module. These maps are natural under restriction, since
both paths send \(da\) to the differential of the same
restricted regular function. Transporting the scalar actions
then gives the asserted natural module-presheaf isomorphism.
\end{proof}

\begin{theorem}[Actual sheafification commutes with open-immersion restriction]
\label{sc0-canonical-port:geom:sheafification-open-immersion-restriction}
\uses{sc0-canonical-port:geom:sheafification-open-restriction,sc0-canonical-port:geom:open-immersion-restricted-sheafification-unit}
\lean{EnvelopeProperty.CanonicalGeometry.openImmersionSheafificationMap_becomes_iso,EnvelopeProperty.CanonicalGeometry.schemeSheafificationCounitIso,EnvelopeProperty.CanonicalGeometry.openImmersionSheafificationIso}
\leanok
For an actual open immersion \(g:U\to X\) and an actual
structure-module presheaf \(P\), the restriction to \(U\)
of the sheafification of \(P\) is naturally isomorphic
to the sheafification of the restricted presheaf.
\end{theorem}
\begin{proof}
Identify the open immersion with the restriction to its actual
open image, using its actual scheme isomorphism onto that image.
Under this identification its open-site functor is the ordinary
inclusion of the opens contained in that image. Restrict the
sheafification unit to this open site. Any local lift and any
local equality required by local bijectivity may be intersected
with the image open, so the restricted unit remains locally
bijective. Its induced map between actual sheafifications is
therefore an isomorphism: locally a section has a lift, and
local equalities give injectivity by sheaf separation.
The actual restricted sheaf is already a sheaf, so its
sheafification counit is an isomorphism. Composing these two
isomorphisms gives the claimed comparison. Transport back
through the actual isomorphism onto the image; this retains
the comparison's unit identity and naturality.
\end{proof}

\begin{theorem}[Actual cotangent sheaves restrict along open immersions]
\label{sc0-canonical-port:geom:cotangent-sheaf-open-immersion-restriction}
\uses{sc0-canonical-port:geom:cotangent-presheaf-open-immersion-restriction,sc0-canonical-port:geom:sheafification-open-immersion-restriction}
\lean{EnvelopeProperty.CanonicalGeometry.cotangentSheafOpenImmersionIso}
\leanok
There is an actual natural structure-module sheaf isomorphism
\[
 g^*(\Omega_{X/K})\simeq\Omega_{U/K}
\]
for restriction along an actual open immersion \(g:U\to X\).
\end{theorem}
\begin{proof}
The preceding sheafification comparison identifies the left
side with the sheafification of the restricted actual cotangent
presheaf. The actual cotangent-presheaf restriction isomorphism
identifies this presheaf with the actual cotangent presheaf of
\(f\circ g\). Applying sheafification gives the displayed
actual sheaf isomorphism.
\end{proof}

\begin{theorem}[An actual affine-open structure is its coordinate structure]
\label{sc0-canonical-port:geom:affine-chart-structural-constants}

\lean{EnvelopeProperty.CanonicalGeometry.affineChartStructuralMap,EnvelopeProperty.CanonicalGeometry.affineChartAlgebra,EnvelopeProperty.CanonicalGeometry.affineChartStructuralMap_eq_affineStructuralMap}
\leanok
For a genuine affine open \(V\subset X\), its actual
chart morphism \(\operatorname{Spec}\Gamma(X,V)\to X\),
followed by \(f\), is the spectrum morphism of the actual
structural map \(K\to\Gamma(X,V)\).
\end{theorem}
\begin{proof}
The spectrum/global-section equivalence identifies morphisms
from the actual affine chart to \(\operatorname{Spec}K\)
with actual ring maps \(K\to\Gamma(X,V)\). The ring map
of the displayed composite is precisely the actual structural
global-section map followed by restriction to \(V\), by
naturality of the actual structure-sheaf maps. This is the
asserted structural constant map. Full faithfulness of spectrum
makes the two actual scheme morphisms equal.
\end{proof}

\begin{theorem}[Smooth scheme morphisms yield actual standard-smooth coordinate charts]
\label{sc0-canonical-port:geom:smooth-scheme-standard-smooth-coordinate-cover}
\uses{sc0-canonical-port:geom:affine-chart-structural-constants,sc0-canonical-port:geom:smooth-scheme-chart-frames}
\lean{EnvelopeProperty.CanonicalGeometry.smooth_cotangent_affine_chart,EnvelopeProperty.CanonicalGeometry.smoothCotangentChart,EnvelopeProperty.CanonicalGeometry.smoothCotangentCharts_coversTop}
\leanok
If \(f\) is smooth of relative dimension \(n\), every
point of \(X\) lies in an actual affine open chart whose
actual coordinate algebra over \(K\) is standard smooth
of relative dimension \(n\). These charts cover \(X\).
\end{theorem}
\begin{proof}
The actual smooth-relative-dimension chart property supplies
affine source and target opens around the given point and its
image, with the actual induced coordinate-ring map standard
smooth of relative dimension \(n\). The target open is
nonempty. Since the spectrum of the field \(K\) has one
point, that target open is all of \(\operatorname{Spec}K\).
Transport its global-section ring through the actual
\(\Gamma(\operatorname{Spec}K,\mathrm{top})\simeq K\)
isomorphism. The preceding structural-constant equality
identifies the transported map with the actual chart's
\(K\)-algebra structure. Transporting a standard-smooth
presentation along the ring isomorphism retains its equations,
Jacobian unit, and relative dimension. This yields the
asserted actual standard-smooth coordinate chart at every point.
\end{proof}

\begin{lemma}[Actual free sheaves retain their basis under open restriction]
\label{sc0-canonical-port:geom:actual-free-sheaf-open-transport}
\uses{sc0-canonical-port:geom:open-immersion-module-presheaf-restriction}
\lean{EnvelopeProperty.CanonicalGeometry.openImmersionRestrictFreeIso,EnvelopeProperty.CanonicalGeometry.schemeOverEquivInverseFreeIso}
\leanok
For an actual open immersion, restriction of the actual
free structure-module sheaf on an index set \(I\)
is the actual free sheaf on the same \(I\).
Likewise, the actual equivalence between module sheaves on
an open subscheme and module sheaves on its open site
transports the free sheaf to the free sheaf on \(I\).
\end{lemma}
\begin{proof}
A free sheaf is the direct sum indexed by \(I\) of
copies of the actual structure sheaf, with its native
unit basis sections. Open restriction sends the structure
sheaf to the structure sheaf by the actual coefficient
isomorphisms. It preserves these sums: the restricted maps
from every summand have the same universal property, since
restriction to an open has its native adjoint and preserves
colimits. Thus the basis-preserving map is an isomorphism.
For the open-site equivalence, the actual unit isomorphism
identifies restriction followed by its inverse with the
identity. Apply that identity and the just obtained free
restriction isomorphism to identify the inverse image of
the free sheaf with the free sheaf on the same basis set.
\end{proof}

\begin{theorem}[The actual cotangent sheaf on an affine chart is its coordinate differential sheaf]
\label{sc0-canonical-port:geom:actual-affine-chart-cotangent-comparison}
\uses{sc0-canonical-port:geom:cotangent-sheaf-open-immersion-restriction,sc0-canonical-port:geom:affine-chart-structural-constants,sc0-canonical-port:geom:actual-affine-cotangent-sheaf-comparison}
\lean{EnvelopeProperty.CanonicalGeometry.affineChartCotangentSheafIso}
\leanok
For any actual affine open \(V\subset X\), give
\(\Gamma(X,V)\) the actual structural \(K\)-algebra.
Restriction of the actual cotangent sheaf along the
native chart \(\operatorname{Spec}\Gamma(X,V)\to X\)
is isomorphic to the associated sheaf of the actual
\(\Omega_{\Gamma(X,V)/K}\) module.
No smoothness hypothesis is needed for this comparison.
\end{theorem}
\begin{proof}
Actual cotangent open restriction identifies the left side
with the cotangent sheaf of the inherited affine structural
morphism. The actual affine-chart constant equality makes
that structural morphism the spectrum of the actual
\(K\to\Gamma(X,V)\) map. The proved actual affine
cotangent comparison then identifies its cotangent sheaf
with the associated sheaf of its coordinate differential
module. Compose these genuine sheaf isomorphisms.
\end{proof}

\begin{theorem}[Actual smooth affine charts give rank-n cotangent sheaf frames]
\label{sc0-canonical-port:geom:smooth-scheme-actual-cotangent-chart-frame}
\uses{sc0-canonical-port:geom:smooth-scheme-standard-smooth-coordinate-cover,sc0-canonical-port:geom:cotangent-sheaf-open-immersion-restriction,sc0-canonical-port:geom:standard-smooth-actual-cotangent-sheaf-frame,sc0-canonical-port:geom:affine-chart-structural-constants,sc0-canonical-port:geom:actual-free-sheaf-open-transport,sc0-canonical-port:geom:actual-affine-chart-cotangent-comparison}
\lean{EnvelopeProperty.CanonicalGeometry.standardSmoothCotangentAffineChartFrame,EnvelopeProperty.CanonicalGeometry.smooth_cotangent_local_frame,EnvelopeProperty.CanonicalGeometry.smoothCotangentChartFrame}
\leanok
On each genuine affine chart in this cover, the restriction
of the actual cotangent sheaf \(\Omega_{X/K}\) is
isomorphic to the free structure-module sheaf of rank \(n\).
\end{theorem}
\begin{proof}
The actual open-immersion cotangent restriction isomorphism
identifies the restricted sheaf with the actual cotangent
sheaf of that affine chart's inherited structural morphism.
The structural-constant equality identifies this morphism
with the actual affine coordinate structure. The chart
algebra is actually standard smooth of relative dimension
\(n\), so Theorem~\ref{sc0-canonical-port:geom:standard-smooth-actual-cotangent-sheaf-frame}
supplies its actual free rank-\(n\) cotangent sheaf frame.
Compose these actual isomorphisms to obtain the required
frame on the restriction of \(\Omega_{X/K}\).
\end{proof}

\begin{corollary}[The actual cotangent sheaf of a smooth scheme is locally free]
\label{sc0-canonical-port:geom:smooth-scheme-actual-cotangent-locally-free}
\uses{sc0-canonical-port:geom:smooth-scheme-actual-cotangent-chart-frame,sc0-canonical-port:geom:smooth-scheme-standard-smooth-coordinate-cover}
\lean{EnvelopeProperty.CanonicalGeometry.smoothCotangentLocalGeneratorsData,EnvelopeProperty.CanonicalGeometry.smoothCotangentLocalGeneratorsData_isLocallyFreeData,EnvelopeProperty.CanonicalGeometry.smooth_cotangentSheaf_locallyFree}
\leanok
For an actual smooth relative-dimension-\(n\) structure
morphism over \(K\), the actual cotangent sheaf is locally
free, with actual rank-\(n\) frames on the cover above.
\end{corollary}
\begin{proof}
Each point belongs to one of the produced actual affine
opens. The actual free-frame isomorphism there supplies
\(n\) local basis sections, and the inverse frame
expresses every restricted section as their actual
regular-function linear combination. These give local
generators together with their linear independence at each
point. They are exactly local free presentations of the
actual sheaf on a cover, which proves local freeness.
\end{proof}

\begin{definition}[Valuation threshold subspace]
\label{sc0-canonical-port:fil:valuation-filtration}
\lean{EnvelopeProperty.CanonicalGeometry.valuationFiltration}
\leanok
Let \(K\subseteq F\) be fields, let \(V\) be a vector space over
\(K\), and let \(\iota:V\to F\) be a \(K\)-linear map. Let
\(v:F\to\R\cup\{+\infty\}\) be an additive valuation, with
\(v(0)=+\infty\), \(v(1)=0\),
\[
 v(xy)=v(x)+v(y),\qquad
 v(x+y)\geq\min\{v(x),v(y)\},
\]
and suppose that \(v(c)=0\) for every \(c\in K^*\). For each real
\(a\), define
\[
 \mathcal F^a_v(V,\iota)=\{s\in V:a\leq v(\iota(s))\}.
\]
The valuation has finite value on nonzero elements of \(F\).
\end{definition}

\begin{lemma}[Subspace closure and decreasing thresholds]
\label{sc0-canonical-port:fil:valuation-subspace}
\uses{sc0-canonical-port:fil:valuation-filtration}
\lean{EnvelopeProperty.CanonicalGeometry.mem_valuationFiltration,EnvelopeProperty.CanonicalGeometry.valuationFiltration_antitone,EnvelopeProperty.CanonicalGeometry.valuation_smul_of_ne_zero}
\leanok
Every \(\mathcal F^a_v(V,\iota)\) is a \(K\)-linear subspace. If
\(a\leq b\), then
\(\mathcal F^b_v(V,\iota)\subseteq\mathcal F^a_v(V,\iota)\).
\end{lemma}
\begin{proof}
The zero vector belongs because its representative has value
\(+\infty\). If \(s,t\) have value at least \(a\), linearity and the
valuation inequality give
\(v(\iota(s+t))\geq\min\{v(\iota(s)),v(\iota(t))\}\geq a\).
For \(c=0\), the vector \(cs\) is zero. For \(c\ne0\), linearity,
multiplicativity, and triviality on \(K^*\) give
\(v(\iota(cs))=v(c)+v(\iota(s))=v(\iota(s))\).
These observations prove scalar closure. Finally, \(b\leq v(\iota(s))\)
and \(a\leq b\) imply \(a\leq v(\iota(s))\), proving the inclusion.
\end{proof}

\begin{lemma}[Multiplicative threshold inequality]
\label{sc0-canonical-port:fil:valuation-product}
\uses{sc0-canonical-port:fil:valuation-filtration}
\lean{EnvelopeProperty.CanonicalGeometry.valuation_product_threshold}
\leanok
For \(s,t\in F\) and real \(a,b\), if \(v(s)\geq a\) and
\(v(t)\geq b\), then \(v(st)\geq a+b\).
Consequently, whenever section-representative maps respect products,
the product of sections in thresholds \(a\) and \(b\) belongs to
threshold \(a+b\) in the product degree.
\end{lemma}
\begin{proof}
The valuation identity gives \(v(st)=v(s)+v(t)\). Addition preserves
the order on \(\R\cup\{+\infty\}\), so this is at least \(a+b\).
Applying this assertion to the representatives of a pair of sections
gives the final assertion.
\end{proof}

\begin{lemma}[Finite value of a nonzero represented vector]
\label{sc0-canonical-port:fil:valuation-finite}
\uses{sc0-canonical-port:fil:valuation-filtration}
\lean{EnvelopeProperty.CanonicalGeometry.valuation_ne_top_of_injective}
\leanok
If \(\iota\) is injective and \(s\ne0\), then
\(v(\iota(s))\ne+\infty\).
\end{lemma}
\begin{proof}
Injectivity and \(\iota(0)=0\) imply \(\iota(s)\ne0\). A valuation
on a field has infinite value only at zero, proving the assertion.
\end{proof}

\begin{lemma}[Units have value zero at a center]
\label{sc0-canonical-port:fil:centered-unit-zero}
\lean{EnvelopeProperty.CanonicalGeometry.valuation_unit_eq_zero}
\leanok
Let \(A\) be a commutative ring, let \(f:A\to F\) be a unital ring
homomorphism to a field, and let \(v\) be an additive valuation of
\(F\). Suppose that \(v(f(r))\geq0\) for all \(r\in A\).
Then \(v(f(u))=0\) for every unit \(u\in A^\times\).
In particular, this applies to the local ring at the center of a
valuation.
\end{lemma}
\begin{proof}
The images of \(u\) and \(u^{-1}\) multiply to one, so
\(v(f(u))+v(f(u^{-1}))=v(1)=0\). Both summands are nonnegative by
the hypothesis. Thus each is zero.
\end{proof}

\begin{lemma}[Independence of a local frame]
\label{sc0-canonical-port:fil:frame-invariance}
\uses{sc0-canonical-port:fil:centered-unit-zero}
\lean{EnvelopeProperty.CanonicalGeometry.valuation_frame_independent}
\leanok
Under the hypotheses of Lemma~\ref{sc0-canonical-port:fil:centered-unit-zero}, for every
unit \(u\in A^\times\) and every \(s\in F\),
\[
 v(f(u)s)=v(s).
\]
Thus two local representatives differing by a unit in the local ring
have the same value.
\end{lemma}
\begin{proof}
Multiplicativity and Lemma~\ref{sc0-canonical-port:fil:centered-unit-zero} give
\(v(f(u)s)=v(f(u))+v(s)=0+v(s)=v(s)\). This includes \(s=0\).
\end{proof}

\begin{lemma}[The unique least-value term survives a finite sum]
\label{sc0-canonical-port:fil:unique-minimum-no-cancellation}
\uses{sc0-canonical-port:fil:valuation-filtration}
\lean{EnvelopeProperty.CanonicalGeometry.valuation_sum_eq_of_unique_min}
\leanok
Let \(f_j\in F\) be a finite family indexed by a finite set \(S\).
If \(i\in S\), \(v(f_i)\) is finite, and
\(v(f_i)<v(f_j)\) for every \(j\in S\setminus\{i\}\), then
\[
 v\left(\sum_{j\in S}f_j\right)=v(f_i).
\]
\end{lemma}
\begin{proof}
If \(v(x)<v(y)\), the valuation inequality gives
\(v(x+y)\geq v(x)\). A strict inequality would, by applying the
same inequality to \(x=(x+y)-y\), imply
\(v(x)\geq\min\{v(x+y),v(y)\}>v(x)\), a contradiction.
Thus \(v(x+y)=v(x)\). Repeated application of the valuation
inequality shows that the sum over \(S\setminus\{i\}\) has value
at least the minimum of its term values, which is strictly greater
than \(v(f_i)\). The empty sum has value \(+\infty\), so it
satisfies the same strict bound. Apply the preceding equality with
\(x=f_i\) and \(y\) equal to that sum.
\end{proof}

\begin{lemma}[Distinct values give linearly independent vectors]
\label{sc0-canonical-port:fil:distinct-values-independent}
\uses{sc0-canonical-port:fil:valuation-filtration,sc0-canonical-port:fil:valuation-finite,sc0-canonical-port:fil:unique-minimum-no-cancellation}
\lean{EnvelopeProperty.CanonicalGeometry.linearIndependent_of_injective_valuation}
\leanok
Suppose \(\iota\) is injective. Any family of nonzero vectors
whose representative values are pairwise distinct is linearly
independent over \(K\).
\end{lemma}
\begin{proof}
Consider a linear relation with at least one nonzero coefficient.
Triviality on \(K^*\) preserves the value of each term with nonzero
coefficient. Among those terms, choose the unique one of smallest
value \(a\). Lemma~\ref{sc0-canonical-port:fil:unique-minimum-no-cancellation} says
that the value of the entire sum is \(a\), which is finite. Its
representative is therefore nonzero. Linearity of \(\iota\) shows
that the original linear combination cannot be zero.
\end{proof}

\begin{definition}[Actual valuation value set]
\label{sc0-canonical-port:fil:valuation-value-set}
\uses{sc0-canonical-port:fil:valuation-filtration}
\lean{EnvelopeProperty.CanonicalGeometry.valuationValues}
\leanok
Let
\[
 \mathcal V_v(V,\iota)
 =\{a\in\R:\text{there exists }s\ne0\text{ with }v(\iota(s))=a\}.
\]
For each \(a\) in this set, choose one vector \(s_a\) realizing
its value. This choice uses only membership in the actual value set.
\end{definition}

\begin{lemma}[One vector per actual value is independent]
\label{sc0-canonical-port:fil:value-representatives-independent}
\uses{sc0-canonical-port:fil:valuation-value-set,sc0-canonical-port:fil:distinct-values-independent}
\lean{EnvelopeProperty.CanonicalGeometry.valuationValueRepresentative,EnvelopeProperty.CanonicalGeometry.valuationValueRepresentative_spec,EnvelopeProperty.CanonicalGeometry.valuationValueRepresentative_linearIndependent}
\leanok
If \(\iota\) is injective, the selected family \((s_a)_{a\in\mathcal V_v(V,\iota)}\)
is linearly independent over \(K\).
\end{lemma}
\begin{proof}
Each selected vector is nonzero and has value equal to its index.
Distinct indices therefore give distinct actual values. Every finite
subfamily is independent by Lemma~\ref{sc0-canonical-port:fil:distinct-values-independent}.
Linear independence means exactly that every finite linear relation
has all coefficients zero, so the entire family is independent.
\end{proof}

\begin{lemma}[Finite jumping values]
\label{sc0-canonical-port:fil:finite-jumping-values}
\uses{sc0-canonical-port:fil:valuation-finite,sc0-canonical-port:fil:valuation-value-set,sc0-canonical-port:fil:value-representatives-independent}
\lean{EnvelopeProperty.CanonicalGeometry.valuationValues_finite,EnvelopeProperty.CanonicalGeometry.valuationValues_card_le_finrank}
\leanok
Suppose \(V\) has finite dimension \(N\) and \(\iota\) is injective.
The values of the nonzero represented vectors form a finite set of
at most \(N\) real numbers.
\end{lemma}
\begin{proof}
If more than \(N\) distinct values occurred, selecting one nonzero
vector at each of \(N+1\) such values would contradict
Lemma~\ref{sc0-canonical-port:fil:value-representatives-independent} and the dimension
bound for a linearly independent family. This proves finiteness
and the cardinality bound.
\end{proof}

\begin{lemma}[Actual exterior maps of semilinear maps]
\label{sc0-canonical-port:ext:semilinear-exterior-map}
\lean{EnvelopeProperty.CanonicalGeometry.exteriorSemilinearAlternating,EnvelopeProperty.CanonicalGeometry.exteriorSemilinearMap,EnvelopeProperty.CanonicalGeometry.exteriorSemilinearMap_ιMulti}
\leanok
Let \(\rho:R\to S\) be a homomorphism of commutative rings
and \(h:M\to N\) a \(\rho\)-semilinear map of actual
modules. For every \(d\geq0\), there is an actual
\(\rho\)-semilinear map
\[
 \bigwedge_R^d M\longrightarrow\bigwedge_S^d N,
 \qquad m_1\wedge\cdots\wedge m_d
 \longmapsto h(m_1)\wedge\cdots\wedge h(m_d).
\]
These maps preserve identity maps and composition.
\end{lemma}
\begin{proof}
Restrict the target scalar action through \(\rho\).
The displayed map of \(d\)-tuples is then
\(R\)-multilinear and alternating. The actual exterior-power
universal property supplies the unique linear map on the
source exterior power, which is exactly the required
semilinear map. Pure wedges span the exterior power,
so equality on those wedges proves the identity and
composition assertions. For \(d=0\), the map is
\(\rho\) on the zeroth exterior powers, with the
same universal property and inverse-free functorial identities.
\end{proof}

\begin{lemma}[An actual linear equivalence gives an actual exterior equivalence]
\label{sc0-canonical-port:ext:linear-exterior-equivalence}
\uses{sc0-canonical-port:ext:semilinear-exterior-map}
\lean{EnvelopeProperty.CanonicalGeometry.exteriorLinearEquiv,EnvelopeProperty.CanonicalGeometry.exteriorLinearEquiv_ιMulti}
\leanok
An actual \(R\)-linear equivalence \(M\simeq N\)
induces an actual linear equivalence of each exterior
power, taking a pure wedge to the wedge of its images.
\end{lemma}
\begin{proof}
Apply the preceding exterior map to the actual linear
equivalence and its inverse. Their composites are the
exterior maps of the identity, as can be checked on
every generating pure wedge. Thus the two exterior
maps are inverse and have the stated wedge formula.
\end{proof}

\begin{theorem}[Actual exterior base change of a free module]
\label{sc0-canonical-port:ext:finite-free-exterior-base-change}
\lean{EnvelopeProperty.CanonicalGeometry.exteriorBasisBaseChangeEquiv,EnvelopeProperty.CanonicalGeometry.exteriorBasisBaseChangeEquiv_basis,EnvelopeProperty.CanonicalGeometry.exteriorBasisBaseChangeEquiv_wedge}
\leanok
Let \(R\to S\) be an actual ring map and \(M\) an actual
\(R\)-module with an ordered basis. For every \(d\geq0\),
the actual canonical comparison
\[
 S\otimes_R\bigwedge_R^d M
 \simeq\bigwedge_S^d(S\otimes_R M)
\]
sends \(1\otimes(m_1\wedge\cdots\wedge m_d)\) to
\((1\otimes m_1)\wedge\cdots\wedge(1\otimes m_d)\).
It takes every scalar-extended basis wedge to the
corresponding actual scalar-extended wedge basis vector.
\end{theorem}
\begin{proof}
The increasing \(d\)-fold wedges of the ordered basis
form an actual basis of the source exterior module.
Expansion of each vector has finite support; multilinear
expansion proves wedge spanning, repeated indices vanish,
and each alternating coordinate determinant extracts
the coefficient of the corresponding increasing wedge,
proving independence. The base-changed module has the
actual base-changed basis, so the same argument gives
its exterior basis. Mapping the scalar-extended source
wedge basis to that target wedge basis constructs an
actual linear equivalence. For an arbitrary pure wedge,
the coefficients in these two bases are the coordinate
determinants. The ring map preserves every determinant
because it preserves finite sums, products, and signs.
Thus the equivalence has the stated formula on every
pure wedge. Scalar linearity extends this to all
actual tensors. For \(d=0\) this is the ordinary
scalar-extension identification with \(S\).
\end{proof}

\begin{theorem}[Actual localization of a free module localizes its exterior modules]
\label{sc0-canonical-port:ext:free-module-exterior-localization}
\uses{sc0-canonical-port:ext:semilinear-exterior-map,sc0-canonical-port:ext:linear-exterior-equivalence,sc0-canonical-port:ext:finite-free-exterior-base-change}
\lean{EnvelopeProperty.CanonicalGeometry.algebraExteriorLinearMap,EnvelopeProperty.CanonicalGeometry.algebraExteriorLinearMap_wedge,EnvelopeProperty.CanonicalGeometry.exteriorLocalizedModuleEquiv,EnvelopeProperty.CanonicalGeometry.exteriorLocalizedModuleEquiv_wedge,EnvelopeProperty.CanonicalGeometry.algebraExteriorLinearMap_isLocalizedModule}
\leanok
Let \(S\subset R\) be a multiplicative subset and
\(B=S^{-1}R\) an actual localization. Let \(M\)
have an actual ordered \(R\)-basis, and let
\(f:M\to N\) be an actual \(R\)-linear localization
map, with the compatible actual \(B\)-module action
on \(N\). Its literal wedge map takes values in
\(\bigwedge_B^dN\) and is itself an actual module
localization map. There is an actual \(B\)-linear
comparison
\[
 B\otimes_R\bigwedge_R^dM\simeq\bigwedge_B^dN,
 \quad 1\otimes\bigwedge_jm_j\longmapsto\bigwedge_jf(m_j).
\]
\end{theorem}
\begin{proof}
The actual localization property of \(f\) identifies
\(N\) with \(B\otimes_RM\), carrying \(f(m)\)
to \(1\otimes m\). Compose the proved free exterior
base-change comparison with the exterior equivalence
of this actual module equivalence. Its wedge formula
is exactly the displayed formula. In particular the
actual wedge map from \(\bigwedge_R^dM\) agrees,
on every generating wedge, with the scalar-extension
unit followed by this equivalence. These maps therefore
agree everywhere. Scalar extension by the actual
localization has its actual localization universal
property; composing with this proved target equivalence
retains that property. Hence the literal wedge map,
not only a chosen basis map, is a module localization map.
\end{proof}

\begin{corollary}[Actual exterior differential localization on free charts]
\label{sc0-canonical-port:ext:free-differential-principal-open}
\uses{sc0-canonical-port:ext:finite-free-exterior-base-change,sc0-canonical-port:ext:linear-exterior-equivalence}
\lean{EnvelopeProperty.CanonicalGeometry.exteriorKaehlerLocalizationEquiv,EnvelopeProperty.CanonicalGeometry.exteriorKaehlerLocalizationEquiv_wedge}
\leanok
Let \(K\to A\to B\) be compatible actual maps of
commutative rings, with \(B\) an actual localization
of \(A\) at a multiplicative subset. If \(\Omega_{A/K}\)
has an actual ordered \(A\)-basis, there is an actual
\(B\)-linear isomorphism
\[
 B\otimes_A\bigwedge_A^d\Omega_{A/K}
 \simeq\bigwedge_B^d\Omega_{B/K}.
\]
It sends a scalar-extended pure wedge to the wedge
of its actual localized differential images.
\end{corollary}
\begin{proof}
Actual differential localization gives
\(B\otimes_A\Omega_{A/K}\simeq\Omega_{B/K}\),
with \(1\otimes da\mapsto d(a|_B)\): a derivation
extends to fractions by differentiating their actual
inverse relation, which gives this differential
universal property and inverse maps. Apply the proved
free exterior base-change comparison to the actual
basis of \(\Omega_{A/K}\), then take the exterior
equivalence of this actual differential localization
isomorphism. The two wedge formulas compose to the
asserted literal localized-differential formula.
\end{proof}

\begin{definition}[The exterior presheaf of an actual module presheaf]
\label{sc0-canonical-port:ext:exterior-presheaf}
\uses{sc0-canonical-port:ext:semilinear-exterior-map,sc0-canonical-port:geom:module-sheaf-sections}
\lean{EnvelopeProperty.CanonicalGeometry.exteriorPresheafRestriction,EnvelopeProperty.CanonicalGeometry.exteriorPresheafRestriction_mk,EnvelopeProperty.CanonicalGeometry.exteriorPresheaf,EnvelopeProperty.CanonicalGeometry.exteriorPresheaf_map_mk}
\leanok
For an actual structure-module presheaf \(P\), set
\(P_d(U)=\bigwedge_{\mathcal O_X(U)}^dP(U)\).
On \(V\subset U\), use the actual exterior map of
the semilinear presheaf restriction. Thus restriction
of a pure wedge is the actual wedge of its restricted
sections. Identity and composition, checked on the
generating wedges, make this an actual module presheaf.
\end{definition}

\begin{theorem}[Presheaf maps are compatible actual alternating families]
\label{sc0-canonical-port:ext:exterior-presheaf-universal}
\uses{sc0-canonical-port:ext:exterior-presheaf}
\lean{EnvelopeProperty.CanonicalGeometry.PresheafAlternatingFamily,EnvelopeProperty.CanonicalGeometry.exteriorPresheafDescApp,EnvelopeProperty.CanonicalGeometry.exteriorPresheafDescApp_mk,EnvelopeProperty.CanonicalGeometry.exteriorPresheafDesc,EnvelopeProperty.CanonicalGeometry.exteriorPresheafHomEquiv}
\leanok
For any actual structure-module presheaf \(N\), maps
\(P_d\to N\) are in bijection with families of actual
alternating \(\mathcal O_X(U)\)-multilinear maps
\(P(U)^d\to N(U)\) compatible with actual restriction.
On each open the inverse map is the actual universal
linear map, taking a pure wedge to that alternating value.
\end{theorem}
\begin{proof}
Compose a presheaf map with the actual pure-wedge maps
to obtain the compatible family. Conversely the actual
exterior universal property on each open gives a unique
linear map taking wedges to the prescribed values.
On a smaller open, its two restriction composites
agree on wedges by the family's compatibility.
Wedges span, so the composites agree everywhere,
proving actual presheaf naturality. The objectwise
universal inverse identities show these constructions
are inverse, establishing the claimed bijection.
\end{proof}

\begin{definition}[The actual exterior module sheaf]
\label{sc0-canonical-port:ext:exterior-sheaf}
\uses{sc0-canonical-port:ext:exterior-presheaf}
\lean{EnvelopeProperty.CanonicalGeometry.exteriorSheaf,EnvelopeProperty.CanonicalGeometry.exteriorSheafificationMap,EnvelopeProperty.CanonicalGeometry.exteriorSheafSection,EnvelopeProperty.CanonicalGeometry.exteriorSheafSection_restrict}
\leanok
Define \(\mathcal E_d(P)\) to be the actual module
sheafification of \(P_d\), with its actual canonical
unit map. Composing the actual pure-wedge maps with
this unit gives alternating operations on actual
presheaf sections with values in \(\mathcal E_d(P)\).
They commute with actual restriction, by naturality
of the unit and the wedge restriction formula.
\end{definition}

\begin{theorem}[The exterior sheaf represents compatible alternating families]
\label{sc0-canonical-port:ext:exterior-sheaf-universal}
\uses{sc0-canonical-port:ext:exterior-sheaf,sc0-canonical-port:ext:exterior-presheaf-universal}
\lean{EnvelopeProperty.CanonicalGeometry.exteriorSheafHomEquiv}
\leanok
For any actual module sheaf \(N\), morphisms
\(\mathcal E_d(P)\to N\) are in bijection with the
restriction-compatible alternating families of
Theorem~\ref{sc0-canonical-port:ext:exterior-presheaf-universal}.
\end{theorem}
\begin{proof}
Module sheafification identifies morphisms out of its
actual sheafification with presheaf maps from its defining
presheaf to the underlying presheaf of \(N\).
Compose this actual adjunction bijection with the
proved exterior-presheaf bijection. Both are defined
by composition with the canonical unit and wedge
maps, giving the asserted alternating interpretation.
\end{proof}

\begin{theorem}[Actual exterior presheaves and sheaves commute with open-site restriction]
\label{sc0-canonical-port:ext:exterior-open-restriction}
\uses{sc0-canonical-port:ext:exterior-presheaf,sc0-canonical-port:ext:exterior-sheaf,sc0-canonical-port:geom:sheafification-open-restriction}
\lean{EnvelopeProperty.CanonicalGeometry.exteriorPresheafOverIso,EnvelopeProperty.CanonicalGeometry.exteriorSheafOverIso}
\leanok
On the actual site of opens contained in \(U\),
restriction of the actual exterior presheaf is
isomorphic to the exterior presheaf of the restricted
presheaf. Restriction of \(\mathcal E_d(P)\) is
isomorphic to the actual sheafification of that
restricted exterior presheaf.
\end{theorem}
\begin{proof}
On an open contained in \(U\), the coefficient ring
and section module are exactly the original ones
on that same open. Thus the exterior objects coincide;
their restriction maps coincide on every pure wedge,
hence everywhere. These objectwise identity maps
are therefore an actual presheaf isomorphism.
Compose its sheafification with the proved actual
open-site sheafification comparison to obtain the
actual exterior-sheaf restriction isomorphism.
\end{proof}

\begin{definition}[The actual relative exterior cotangent presheaf and sheaf]
\label{sc0-canonical-port:ext:relative-exterior-cotangent-sheaf}
\uses{sc0-canonical-port:geom:cotangent-presheaf,sc0-canonical-port:ext:exterior-presheaf,sc0-canonical-port:ext:exterior-sheaf}
\lean{EnvelopeProperty.CanonicalGeometry.exteriorCotangentPresheaf,EnvelopeProperty.CanonicalGeometry.exteriorCotangentSheaf}
\leanok
For an actual \(f:X\to\operatorname{Spec}K\), with
\(K\) a field, take the objectwise actual exterior
powers of its actual cotangent presheaf and then
sheafify. Denote the resulting actual sheaf by
\(\Omega_{X/K}^{[d]}\). This construction is made
from the actual regular-function differential modules.
\end{definition}

\begin{lemma}[Actual affine coordinate differential maps]
\label{sc0-canonical-port:ext:affine-ordinary-differential-coordinate-map}
\uses{sc0-canonical-port:geom:actual-affine-cotangent-objects,sc0-canonical-port:geom:actual-affine-differential-localized-maps}
\lean{EnvelopeProperty.CanonicalGeometry.affineOrdinaryCotangentToOpen,EnvelopeProperty.CanonicalGeometry.affineOrdinaryCotangentToOpen_naturality,EnvelopeProperty.CanonicalGeometry.affineOrdinaryCotangentToOpen_isLocalized,EnvelopeProperty.CanonicalGeometry.affineOrdinaryCotangentSemilinear}
\leanok
Let \(K\) be a field and \(A\) an actual commutative
\(K\)-algebra. For every open of \(\operatorname{Spec}A\),
there is an actual coordinate map from ordinary
\(\Omega_{A/K}\) to the actual cotangent-presheaf
module on that open. It is semilinear for the actual
coordinate regular-function map, is natural under
restriction, and on \(D(r)\) is an actual module
localization map at the powers of \(r\).
\end{lemma}
\begin{proof}
Transport the proved actual affine cotangent coordinate
map through the actual coordinate differential-module
isomorphism to ordinary \(\Omega_{A/K}\).
This isomorphism preserves the actual differential
images, so the map retains both restriction naturality
and the coordinate scalar action. On a principal
open, composing a genuine module localization map
with this proved source linear equivalence retains
its universal property. This gives the asserted
actual localization map and semilinear description.
\end{proof}

\begin{lemma}[Actual affine coordinate exterior maps and their localization property]
\label{sc0-canonical-port:ext:affine-exterior-coordinate-map}
\uses{sc0-canonical-port:ext:affine-ordinary-differential-coordinate-map,sc0-canonical-port:ext:semilinear-exterior-map,sc0-canonical-port:ext:free-module-exterior-localization,sc0-canonical-port:ext:relative-exterior-cotangent-sheaf}
\lean{EnvelopeProperty.CanonicalGeometry.affineExteriorCotangentModule,EnvelopeProperty.CanonicalGeometry.affineExteriorCotangentToOpen,EnvelopeProperty.CanonicalGeometry.affineExteriorCotangentToOpen_wedge,EnvelopeProperty.CanonicalGeometry.affineExteriorCotangentToOpen_naturality,EnvelopeProperty.CanonicalGeometry.affineExteriorCotangentToOpen_isLocalized}
\leanok
The actual coordinate module is \(\bigwedge_A^d\Omega_{A/K}\).
Its actual map to the relative exterior presheaf on
an open takes every pure wedge to the wedge of the
actual differential coordinate images. These maps
are natural under restriction. If \(\Omega_{A/K}\)
has an actual ordered basis, the map on every
principal open is an actual module localization map.
\end{lemma}
\begin{proof}
Apply the actual semilinear exterior construction
to the preceding actual differential coordinate map.
The pure-wedge formula follows from its universal
construction. On restriction, both composites have
the same pure-wedge values by differential naturality;
wedge spanning proves equality on the entire module.
On a principal open the differential coordinate map
is an actual localization map, its coefficient ring
is the actual localization of \(A\), and the source
has the assumed actual basis. The proved free-module
exterior localization theorem therefore gives the
actual localization property for this literal map.
\end{proof}

\begin{theorem}[Actual affine relative exterior sheaves are associated exterior differential modules]
\label{sc0-canonical-port:ext:actual-affine-exterior-sheaf-comparison}
\uses{sc0-canonical-port:ext:affine-exterior-coordinate-map,sc0-canonical-port:geom:affine-localized-sheafification-isomorphism}
\lean{EnvelopeProperty.CanonicalGeometry.affineExteriorCotangentSheafIso}
\leanok
If \(\Omega_{A/K}\) has an actual ordered basis,
there is an actual structure-module sheaf isomorphism
\[
 \Omega_{\operatorname{Spec}A/K}^{[d]}
 \simeq\widetilde{\bigwedge_A^d\Omega_{A/K}}.
\]
\end{theorem}
\begin{proof}
The preceding coordinate exterior maps are actual
natural maps and supply actual localized-module
comparisons on every principal open. Apply the
proved affine localized-sheafification comparison
with the actual module \(\bigwedge_A^d\Omega_{A/K}\)
and its actual exterior cotangent presheaf.
Its hypotheses have been produced by those actual
maps, so it gives precisely the displayed actual
sheaf isomorphism.
\end{proof}

\begin{theorem}[An actual finite differential basis gives an actual affine top-exterior line frame]
\label{sc0-canonical-port:ext:actual-affine-top-exterior-frame}
\uses{sc0-canonical-port:ext:actual-affine-exterior-sheaf-comparison,sc0-canonical-port:geom:canonical-frame-from-differential-basis,sc0-canonical-port:geom:affine-tilde-line-bundle}
\lean{EnvelopeProperty.CanonicalGeometry.affineTopExteriorSheafFrame,EnvelopeProperty.CanonicalGeometry.affineExteriorCanonicalLineBundle}
\leanok
If \(\Omega_{A/K}\) has an actual basis indexed
by \(\{0,\ldots,n-1\}\), its actual relative top
exterior sheaf has a frame to the unit structure-module
sheaf. This produces an actual line atlas on that sheaf.
\end{theorem}
\begin{proof}
The actual finite differential basis gives an actual
one-element basis of its top exterior module, hence
an actual module trivialization to \(A\).
The associated-sheaf functor transports that actual
trivialization to a frame of the associated top
exterior sheaf. Compose with the preceding actual
relative exterior-sheaf comparison. This gives a
true global affine frame, and its single top open
forms a genuine line atlas.
\end{proof}

\begin{corollary}[Actual standard smoothness produces the affine top-exterior line frame]
\label{sc0-canonical-port:ext:standard-smooth-actual-top-exterior-frame}
\uses{sc0-canonical-port:ext:actual-affine-top-exterior-frame,sc0-canonical-port:geom:standard-smooth-differential-basis}
\lean{EnvelopeProperty.CanonicalGeometry.standardSmoothTopExteriorSheafFrame,EnvelopeProperty.CanonicalGeometry.standardSmoothExteriorCanonicalLineBundle}
\leanok
For an actual standard-smooth relative-dimension-\(n\)
\(K\)-algebra, the actual relative top exterior sheaf
has a produced frame and actual line atlas, without
a differential-basis or line-frame premise.
\end{corollary}
\begin{proof}
The actual standard-smooth presentation supplies the
proved basis indexed by \(\operatorname{Fin}n\).
Apply the preceding actual affine top-exterior frame
construction to that produced basis.
\end{proof}

\begin{definition}[The actual associated affine dual exterior module sheaf]
\label{sc0-canonical-port:ext:actual-affine-dual-exterior-module-sheaf}
\uses{sc0-canonical-port:geom:affine-canonical-module}
\lean{EnvelopeProperty.CanonicalGeometry.affineExteriorAnticanonicalSheaf}
\leanok
Define the actual affine dual exterior module sheaf as
\(\widetilde{\operatorname{Hom}_A(\bigwedge_A^n\Omega_{A/K},A)}\).
This is the associated sheaf of the actual coordinate
module dual. Its identification with an intrinsic
sheaf dual on arbitrary schemes is a separate statement.
\end{definition}

\begin{theorem}[Actual affine dual exterior module sheaves have produced line frames]
\label{sc0-canonical-port:ext:actual-affine-dual-exterior-frame}
\uses{sc0-canonical-port:ext:actual-affine-dual-exterior-module-sheaf,sc0-canonical-port:geom:anticanonical-dual-frame,sc0-canonical-port:geom:affine-tilde-line-bundle,sc0-canonical-port:geom:standard-smooth-differential-basis}
\lean{EnvelopeProperty.CanonicalGeometry.affineExteriorAnticanonicalSheafFrame,EnvelopeProperty.CanonicalGeometry.affineExteriorAnticanonicalLineBundle,EnvelopeProperty.CanonicalGeometry.standardSmoothExteriorAnticanonicalLineBundle}
\leanok
An actual \(\operatorname{Fin}n\)-indexed differential
basis produces an actual frame and line atlas on this
affine associated dual sheaf. Actual standard smoothness
of relative dimension \(n\) produces that basis and
therefore this actual dual line atlas.
\end{theorem}
\begin{proof}
The actual dual of the top-wedge basis trivializes the
coordinate dual module. The associated-sheaf functor
transports its actual linear trivialization to an
actual sheaf frame; taking the whole affine open gives
a line atlas. In the standard-smooth case, its actual
presentation supplies the differential basis, so the
same construction gives the claimed actual atlas.
\end{proof}

\begin{theorem}[Actual affine top-exterior and dual-module frames give a tensor frame]
\label{sc0-canonical-port:ext:actual-affine-exterior-dual-tensor-frame}
\uses{sc0-canonical-port:ext:actual-affine-top-exterior-frame,sc0-canonical-port:ext:actual-affine-dual-exterior-frame,sc0-canonical-port:geom:tensor-sheaf}
\lean{EnvelopeProperty.CanonicalGeometry.affineExteriorDualTensorFrame}
\leanok
Under the actual finite differential-basis hypothesis,
the tensor of the actual affine associated dual exterior
sheaf and the actual relative top exterior sheaf has
an actual frame to the unit structure-module sheaf.
\end{theorem}
\begin{proof}
Tensor their proved actual frames to identify their
actual tensor sheaf with the tensor of two unit sheaves.
The actual tensor unit isomorphism identifies that
sheaf with the unit sheaf. This composition is the
asserted actual tensor frame.
\end{proof}

\begin{theorem}[Actual coordinate dual tensor evaluation is a linear isomorphism]
\label{sc0-canonical-port:ext:actual-affine-coordinate-dual-evaluation}
\uses{sc0-canonical-port:geom:anticanonical-dual-frame,sc0-canonical-port:geom:canonical-frame-from-differential-basis}
\lean{EnvelopeProperty.CanonicalGeometry.affineExteriorEvaluationEquiv,EnvelopeProperty.CanonicalGeometry.affineExteriorEvaluationEquiv_tmul}
\leanok
With \(M=\bigwedge_A^n\Omega_{A/K}\) and an actual
\(\operatorname{Fin}n\)-indexed differential basis,
there is an actual linear isomorphism
\(M^\vee\otimes_AM\simeq A\) whose value at every
pure tensor \(\varphi\otimes x\) is literally
\(\varphi(x)\).
\end{theorem}
\begin{proof}
Let \(g\) be the actual top-wedge generator, and
write \(x=a g\) using its actual one-element basis.
The actual dual trivialization sends \(\varphi\)
to \(\varphi(g)\); the actual primal trivialization
sends \(x\) to \(a\). Tensor these actual linear
isomorphisms and use \(A\otimes_AA\simeq A\).
The resulting map sends \(\varphi\otimes x\)
to \(\varphi(g)a=a\varphi(g)=\varphi(a g)=\varphi(x)\).
Its construction from actual linear isomorphisms
makes it an actual linear isomorphism, proving both
claims.
\end{proof}

\begin{definition}[Actual sections of a module sheaf]
\label{sc0-canonical-port:geom:module-sheaf-sections}
\lean{EnvelopeProperty.CanonicalGeometry.bundleSections,EnvelopeProperty.CanonicalGeometry.bundleSectionsModule,EnvelopeProperty.CanonicalGeometry.bundleSectionsSMulCommClass}
\leanok
Let \(X\) be a scheme and \(\mathcal L\) a sheaf of
\(\mathcal O_X\)-modules. For an open \(U\subseteq X\), let
\(\Gamma(U,\mathcal L)=\mathcal L(U)\), retaining its actual
\(\mathcal O_X(U)\)-module structure and restriction maps.
\end{definition}

\begin{definition}[The actual tensor sheaf]
\label{sc0-canonical-port:geom:tensor-sheaf}
\uses{sc0-canonical-port:geom:module-sheaf-sections}
\lean{EnvelopeProperty.CanonicalGeometry.bundleTensorPresheaf,EnvelopeProperty.CanonicalGeometry.bundleTensor}
\leanok
For sheaves \(\mathcal L,\mathcal N\) of
\(\mathcal O_X\)-modules, form the tensor presheaf
\[
 U\longmapsto\mathcal L(U)\otimes_{\mathcal O_X(U)}\mathcal N(U),
\]
with the maps induced by restriction of sections and scalars.
Define \(\mathcal L\otimes\mathcal N\) to be its sheafification.
\end{definition}

\begin{definition}[Pure tensors of actual local sections]
\label{sc0-canonical-port:geom:pure-tensor-sections}
\uses{sc0-canonical-port:geom:tensor-sheaf,sc0-canonical-port:geom:module-sheaf-sections}
\lean{EnvelopeProperty.CanonicalGeometry.tensorSection}
\leanok
For \(s\in\Gamma(U,\mathcal L)\) and
\(t\in\Gamma(U,\mathcal N)\), write \(s\boxtimes t\) for the
image of \(s\otimes t\) under the canonical sheafification map,
as a section of \(\mathcal L\otimes\mathcal N\) on \(U\).
\end{definition}

\begin{lemma}[Bilinearity of actual tensor sections]
\label{sc0-canonical-port:geom:tensor-section-bilinearity}
\uses{sc0-canonical-port:geom:pure-tensor-sections}
\lean{EnvelopeProperty.CanonicalGeometry.tensorSection_add_left,EnvelopeProperty.CanonicalGeometry.tensorSection_add_right,EnvelopeProperty.CanonicalGeometry.tensorSection_zero_left,EnvelopeProperty.CanonicalGeometry.tensorSection_zero_right,EnvelopeProperty.CanonicalGeometry.tensorSection_smul_left,EnvelopeProperty.CanonicalGeometry.tensorSection_smul_right,EnvelopeProperty.CanonicalGeometry.tensorSectionBilinear}
\leanok
The pure tensor operation is additive in each variable and satisfies
\(0\boxtimes t=0\) and \(s\boxtimes0=0\).
It is also linear in both variables over \(\mathcal O_X(U)\):
\((a s)\boxtimes t=a(s\boxtimes t)=s\boxtimes(a t)\).
Thus it defines an actual bilinear map of the section modules.
\end{lemma}
\begin{proof}
The tensor-product identities give
\((s+s')\otimes t=s\otimes t+s'\otimes t\) and
\(s\otimes(t+t')=s\otimes t+s\otimes t'\), with zero tensors
zero. The sheafification unit is a module homomorphism on each
open, so preserves those sums and zero. Applying it proves the
sum and zero identities. The tensor identities also give
\((a s)\otimes t=a(s\otimes t)=s\otimes(a t)\).
The same module homomorphism preserves scalar multiplication,
proving linearity in both variables and the bilinear-map assertion.
\end{proof}

\begin{lemma}[Tensor section multiplication commutes with restriction]
\label{sc0-canonical-port:geom:tensor-section-restriction}
\uses{sc0-canonical-port:geom:pure-tensor-sections}
\lean{EnvelopeProperty.CanonicalGeometry.tensorSection_restrict}
\leanok
For \(V\subseteq U\),
\[
 (s\boxtimes t)|_V=(s|_V)\boxtimes(t|_V).
\]
\end{lemma}
\begin{proof}
The presheaf tensor restriction sends a pure tensor to the pure
tensor of its restricted factors. Naturality of the sheafification
unit says that applying it before restricting is the same as
restricting the tensor first and then applying the unit on \(V\).
This is exactly the displayed identity.
\end{proof}

\begin{definition}[Tensor sheaf powers and their genuine global sections]
\label{sc0-canonical-port:geom:tensor-sheaf-powers}
\uses{sc0-canonical-port:geom:tensor-sheaf,sc0-canonical-port:geom:module-sheaf-sections}
\lean{EnvelopeProperty.CanonicalGeometry.bundleTensorPower,EnvelopeProperty.CanonicalGeometry.bundlePowerGlobalSections}
\leanok
For an actual module sheaf \(\mathcal L\), set
\[
 \mathcal L^{\otimes0}=\mathcal O_X,\qquad
 \mathcal L^{\otimes(n+1)}=\mathcal L^{\otimes n}\otimes\mathcal L.
\]
Its degree-\(n\) section space is the actual module
\(\Gamma(X,\mathcal L^{\otimes n})\).
No finite-dimensionality, positivity, or Hilbert asymptotic is
included in this construction.
\end{definition}

\begin{definition}[Degree-successor multiplication on actual sections]
\label{sc0-canonical-port:geom:successor-section-product}
\uses{sc0-canonical-port:geom:tensor-sheaf-powers,sc0-canonical-port:geom:pure-tensor-sections}
\lean{EnvelopeProperty.CanonicalGeometry.bundlePowerSectionMul,EnvelopeProperty.CanonicalGeometry.bundlePowerGlobalSectionMul}
\leanok
For \(s\in\Gamma(U,\mathcal L^{\otimes n})\) and
\(t\in\Gamma(U,\mathcal L)\), define their successor product
as \(s\boxtimes t\in\Gamma(U,\mathcal L^{\otimes(n+1)})\).
In particular this gives the actual degree-successor product on
global sections when \(U=X\).
\end{definition}

\begin{lemma}[The actual successor product respects restrictions]
\label{sc0-canonical-port:geom:successor-product-restriction}
\uses{sc0-canonical-port:geom:successor-section-product,sc0-canonical-port:geom:tensor-section-restriction}
\lean{EnvelopeProperty.CanonicalGeometry.bundlePowerSectionMul_restrict}
\leanok
The degree-successor product restricts to the degree-successor
product of the restricted sections on each smaller open.
\end{lemma}
\begin{proof}
Apply Lemma~\ref{sc0-canonical-port:geom:tensor-section-restriction} with
\(\mathcal N=\mathcal L\) and first factor
\(\mathcal L^{\otimes n}\). Its target sheaf is precisely the
recursively defined \(\mathcal L^{\otimes(n+1)}\).
\end{proof}

\begin{definition}[Genuine local rank-one frames]
\label{sc0-canonical-port:geom:local-line-bundle-frames}
\uses{sc0-canonical-port:geom:module-sheaf-sections}
\lean{EnvelopeProperty.CanonicalGeometry.LocalBundleFrame,EnvelopeProperty.CanonicalGeometry.LineBundleAtlas,EnvelopeProperty.CanonicalGeometry.IntrinsicLineBundle}
\leanok
A local frame of \(\mathcal L\) on \(U\) is an isomorphism of
actual sheaves of \(\mathcal O_U\)-modules
\(\mathcal L|_U\simeq\mathcal O_U\).
A rank-one atlas consists of an open cover \((U_i)\) and one such
frame on each \(U_i\). An intrinsic line bundle is an actual
\(\mathcal O_X\)-module sheaf supplied with such an atlas.
\end{definition}

\begin{definition}[The affine canonical and anticanonical modules]
\label{sc0-canonical-port:geom:affine-canonical-module}
\lean{EnvelopeProperty.CanonicalGeometry.affineCanonicalModule,EnvelopeProperty.CanonicalGeometry.affineAnticanonicalModule}
\leanok
Let \(A\) be a commutative \(K\)-algebra, and let
\(\Omega_{A/K}\) be its actual module of K\"ahler differentials.
For an integer \(n\geq0\), define
\[
 \omega_{A/K}^{(n)}=\bigwedge_A^n\Omega_{A/K},\qquad
 (\omega_{A/K}^{(n)})^{-1}
   =\operatorname{Hom}_A(\omega_{A/K}^{(n)},A).
\]
These are the top differential and its module dual when the
relative differential rank is \(n\).
\end{definition}

\begin{theorem}[Actual rank-one canonical frames from a differential basis]
\label{sc0-canonical-port:geom:canonical-frame-from-differential-basis}
\uses{sc0-canonical-port:geom:affine-canonical-module}
\lean{EnvelopeProperty.CanonicalGeometry.affineCanonicalBasis,EnvelopeProperty.CanonicalGeometry.affineCanonicalGenerator,EnvelopeProperty.CanonicalGeometry.affineCanonicalGenerator_eq_wedge,EnvelopeProperty.CanonicalGeometry.affineCanonicalTrivialization,EnvelopeProperty.CanonicalGeometry.affineCanonicalTrivialization_generator,EnvelopeProperty.CanonicalGeometry.affineCanonical_eq_smul_generator}
\leanok
If \((e_i)_{i=0}^{n-1}\) is an \(A\)-basis of
\(\Omega_{A/K}\), then
\(\kappa=e_0\wedge\cdots\wedge e_{n-1}\) is a one-element
basis of \(\omega_{A/K}^{(n)}\). It gives an actual
\(A\)-linear isomorphism \(\omega_{A/K}^{(n)}\simeq A\).
\end{theorem}
\begin{proof}
Exterior products of \(n\) vectors span the exterior power by its
construction. Expand each vector in the specified differential
basis and distribute the wedge. Terms with a repeated basis index
are zero by alternation. Every remaining term uses all \(n\)
basis indices and is its permutation sign times \(\kappa\).
Thus \(\kappa\) spans.
The determinant of the \(n\) coordinate columns is an alternating
\(A\)-multilinear form, so it induces a linear map
\(\delta:\omega_{A/K}^{(n)}\to A\) with \(\delta(\kappa)=1\).
If \(a\kappa=0\), applying \(\delta\) gives \(a=0\), proving
independence. Therefore \(a\mapsto a\kappa\) and \(\delta\)
are inverse linear isomorphisms. For \(n=0\), the zeroth exterior
power is \(A\), the empty wedge is one, and the determinant is
one; the same assertions hold.
\end{proof}

\begin{theorem}[The actual anticanonical dual frame]
\label{sc0-canonical-port:geom:anticanonical-dual-frame}
\uses{sc0-canonical-port:geom:canonical-frame-from-differential-basis}
\lean{EnvelopeProperty.CanonicalGeometry.affineAnticanonicalBasis,EnvelopeProperty.CanonicalGeometry.affineAnticanonicalGenerator,EnvelopeProperty.CanonicalGeometry.affineAnticanonicalTrivialization,EnvelopeProperty.CanonicalGeometry.affineAnticanonicalGenerator_apply_generator,EnvelopeProperty.CanonicalGeometry.affineAnticanonicalTrivialization_apply,EnvelopeProperty.CanonicalGeometry.affineAnticanonical_eq_smul_generator,EnvelopeProperty.CanonicalGeometry.affineCanonical_finite_free,EnvelopeProperty.CanonicalGeometry.affineAnticanonical_finite_free,EnvelopeProperty.CanonicalGeometry.affineCanonical_finrank,EnvelopeProperty.CanonicalGeometry.affineAnticanonical_finrank}
\leanok
Under the differential-basis hypothesis above, the functional
\(\delta\) with \(\delta(\kappa)=1\) is a one-element basis of
\((\omega_{A/K}^{(n)})^{-1}\), giving an actual
\(A\)-linear isomorphism of that dual module with \(A\).
Both modules are free and finite. If \(A\ne0\), each has rank one.
\end{theorem}
\begin{proof}
Every element of the canonical module is uniquely \(a\kappa\).
An \(A\)-linear functional \(\psi\) satisfies
\(\psi(a\kappa)=a\psi(\kappa)\), so
\(\psi=\psi(\kappa)\delta\). If \(b\delta=0\), evaluating
at \(\kappa\) gives \(b=0\). Thus \(\delta\) is a basis and
\(\psi\mapsto\psi(\kappa)\) is the indicated isomorphism.
A one-element basis proves freeness and finite generation, and in
a nonzero coefficient ring its index cardinality is rank one.
\end{proof}

\begin{corollary}[The polynomial canonical wedge and its normalized dual]
\label{sc0-canonical-port:geom:polynomial-canonical-frame}
\uses{sc0-canonical-port:geom:canonical-frame-from-differential-basis,sc0-canonical-port:geom:anticanonical-dual-frame}
\lean{EnvelopeProperty.CanonicalGeometry.polynomialCanonicalBasis,EnvelopeProperty.CanonicalGeometry.polynomialCanonicalGenerator,EnvelopeProperty.CanonicalGeometry.polynomialCanonicalGenerator_eq_wedge,EnvelopeProperty.CanonicalGeometry.polynomialAnticanonicalBasis,EnvelopeProperty.CanonicalGeometry.polynomialAnticanonicalGenerator,EnvelopeProperty.CanonicalGeometry.polynomialCanonical_finite_free,EnvelopeProperty.CanonicalGeometry.polynomialAnticanonical_finite_free,EnvelopeProperty.CanonicalGeometry.polynomialCanonical_finrank,EnvelopeProperty.CanonicalGeometry.polynomialAnticanonical_finrank,EnvelopeProperty.CanonicalGeometry.polynomialAnticanonicalGenerator_apply_generator}
\leanok
For \(A=K[X_0,\ldots,X_{n-1}]\), the canonical module has the
actual frame \(\mathrm dX_0\wedge\cdots\wedge\mathrm dX_{n-1}\),
and its dual has the actual frame that evaluates this wedge to one.
These modules are finite free, and each has rank one if \(K\ne0\).
This includes \(n=0\).
\end{corollary}
\begin{proof}
Apply Theorems~\ref{sc0-canonical-port:geom:canonical-frame-from-differential-basis}
and \ref{sc0-canonical-port:geom:anticanonical-dual-frame} to the differential basis
produced in Theorem~\ref{sc0-canonical-port:geom:polynomial-differential-basis}.
Their empty-basis cases give the asserted zero-dimensional case.
\end{proof}

\begin{lemma}[The unique full exterior index]
\label{sc0-canonical-port:geom:top-exterior-full-index}
\lean{EnvelopeProperty.CanonicalGeometry.topExteriorIndexEquiv,EnvelopeProperty.CanonicalGeometry.topExteriorIndexEquiv_symm_orderEmb}
\leanok
The family of \(n\)-element subsets of \(\{0,\ldots,n-1\}\)
has exactly one element, the full set. Its increasing enumeration
is the identity on that ordered index set, including \(n=0\).
\end{lemma}
\begin{proof}
A proper subset of a finite set has strictly smaller cardinality,
so an \(n\)-element subset of this \(n\)-element set is the
full set. An increasing bijection from this finite ordered set
to itself must fix the least element, then the next least, and
successively every element; it is the identity. For the empty
index set there is one subset of cardinality zero, and one empty
enumeration, giving the same assertions.
\end{proof}

\begin{lemma}[Canonical wedge coordinates are actual determinants]
\label{sc0-canonical-port:geom:canonical-wedge-determinant}
\uses{sc0-canonical-port:geom:canonical-frame-from-differential-basis}
\lean{EnvelopeProperty.CanonicalGeometry.affineCanonicalTrivialization_apply_wedge}
\leanok
For the differential basis \(b=(e_j)\) and any \(n\) differential
vectors \(v_i\), the canonical trivialization sends
\(v_0\wedge\cdots\wedge v_{n-1}\) to
\(\det(b_j^*(v_i))_{i,j}\).
\end{lemma}
\begin{proof}
Expand each \(v_i=\sum_j b_j^*(v_i)e_j\) and distribute the
wedge. Repeated indices vanish, and each permutation of the full
index set contributes its sign times the product of its coordinate
entries. The coefficient of \(\kappa\) is therefore the
permutation expansion of the stated determinant. The trivialization
extracts exactly that coefficient.
\end{proof}

\begin{lemma}[Actual canonical transition determinants are units]
\label{sc0-canonical-port:geom:canonical-transition-unit}
\uses{sc0-canonical-port:geom:canonical-frame-from-differential-basis,sc0-canonical-port:geom:canonical-wedge-determinant}
\lean{EnvelopeProperty.CanonicalGeometry.affineCanonical_transition_isUnit,EnvelopeProperty.CanonicalGeometry.affineCanonical_transition_det_isUnit}
\leanok
For two actual differential bases \(b,b'\), the scalar changing
their canonical generators is a unit, equal to
\(\det(b_j^*(b'_i))_{i,j}\).
\end{lemma}
\begin{proof}
Each canonical generator is a basis, so write
\(\kappa'=a\kappa\) and \(\kappa=d\kappa'\) uniquely.
Substitution gives \(\kappa=da\kappa\), and independence of
\(\kappa\) forces \(da=1\). Thus \(a\) is a unit with inverse
\(d\). Lemma~\ref{sc0-canonical-port:geom:canonical-wedge-determinant} computes
\(a\) as the displayed determinant.
\end{proof}

\begin{theorem}[Differential bases transported through actual formally etale maps]
\label{sc0-canonical-port:geom:etale-differential-basis}
\lean{EnvelopeProperty.CanonicalGeometry.affineDifferentialBasisOfFormallyEtale,EnvelopeProperty.CanonicalGeometry.affineDifferentialBasisOfFormallyEtale_apply}
\leanok
Let \(K\to A\to B\) be compatible homomorphisms of commutative
rings, with \(A\to B\) formally etale. If \((e_i)_{i=0}^{n-1}\)
is an \(A\)-basis of \(\Omega_{A/K}\), the images of these
actual differential vectors form a \(B\)-basis of
\(\Omega_{B/K}\).
\end{theorem}
\begin{proof}
Formal etaleness gives the differential base-change isomorphism
\(B\otimes_A\Omega_{A/K}\simeq\Omega_{B/K}\).
One way to see this is by testing linear maps into any \(B\)-module:
they correspond to \(K\)-derivations, and a derivation of \(A\)
has a unique extension to \(B\) by the unique lifting property
for the square-zero extension \(B\oplus N\). This is precisely
the universal isomorphism of the represented differential modules.
Base change of a free module with basis \((e_i)\) has basis
\((1\otimes e_i)\); the isomorphism sends these vectors to their
actual differential images. Transporting this basis proves the
claim and the stated formula for its vectors.
\end{proof}

\begin{corollary}[Actual differential bases on localization charts]
\label{sc0-canonical-port:geom:localization-differential-basis}
\uses{sc0-canonical-port:geom:etale-differential-basis}
\lean{EnvelopeProperty.CanonicalGeometry.affineDifferentialBasisOfLocalization}
\leanok
If \(B=S^{-1}A\) is an actual localization, an \(A\)-basis of
\(\Omega_{A/K}\) produces a \(B\)-basis of \(\Omega_{B/K}\).
\end{corollary}
\begin{proof}
A derivation extends uniquely through a localization by
\(D(a/s)=D(a)/s-aD(s)/s^2\). Equivalently the localization is
formally etale: units lift uniquely across a square-zero ideal,
so its defining inversion property gives unique infinitesimal lifts.
Apply Theorem~\ref{sc0-canonical-port:geom:etale-differential-basis}.
\end{proof}

\begin{corollary}[Actual formally etale polynomial charts produce canonical frames]
\label{sc0-canonical-port:geom:etale-polynomial-canonical-frame}
\uses{sc0-canonical-port:geom:etale-differential-basis,sc0-canonical-port:geom:canonical-frame-from-differential-basis,sc0-canonical-port:geom:anticanonical-dual-frame}
\lean{EnvelopeProperty.CanonicalGeometry.polynomialEtaleDifferentialBasis,EnvelopeProperty.CanonicalGeometry.polynomialEtaleDifferentialBasis_apply,EnvelopeProperty.CanonicalGeometry.polynomialEtale_canonical_anticanonical}
\leanok
Let \(B\ne0\) be a commutative \(K\)-algebra with a compatible
formally etale map \(K[X_0,\ldots,X_{n-1}]\to B\).
The actual differentials of the images of \(X_i\) form a basis
of \(\Omega_{B/K}\). The canonical module and its actual dual
are finite free of rank one, with the top wedge and its dual frames.
\end{corollary}
\begin{proof}
The polynomial differential basis comes from
Theorem~\ref{sc0-canonical-port:geom:polynomial-differential-basis}.
Theorem~\ref{sc0-canonical-port:geom:etale-differential-basis} transports it to the
actual differentials of the coordinate images in \(B\).
Apply Theorems~\ref{sc0-canonical-port:geom:canonical-frame-from-differential-basis}
and \ref{sc0-canonical-port:geom:anticanonical-dual-frame} to this produced basis.
Their one-element bases give finite freeness and, since \(B\ne0\),
rank one for both modules.
\end{proof}

\begin{corollary}[Actual polynomial localizations produce canonical frames]
\label{sc0-canonical-port:geom:polynomial-localization-canonical-frame}
\uses{sc0-canonical-port:geom:localization-differential-basis,sc0-canonical-port:geom:canonical-frame-from-differential-basis,sc0-canonical-port:geom:anticanonical-dual-frame}
\lean{EnvelopeProperty.CanonicalGeometry.polynomialLocalizationDifferentialBasis,EnvelopeProperty.CanonicalGeometry.polynomialLocalization_canonical_anticanonical}
\leanok
If \(B\ne0\) is an actual localization of
\(K[X_0,\ldots,X_{n-1}]\), its actual differential module has
rank \(n\) with the localized coordinate differential basis.
Its canonical module and its actual dual are finite free of rank one.
\end{corollary}
\begin{proof}
Apply Corollary~\ref{sc0-canonical-port:geom:localization-differential-basis} to the
polynomial basis, and then apply the canonical and dual frame
constructions to the resulting basis over \(B\). Their finite
one-element bases give the asserted finite freeness and rank.
\end{proof}

\begin{lemma}[Tensor sheaves transport actual sheaf isomorphisms]
\label{sc0-canonical-port:geom:tensor-sheaf-isomorphisms}
\uses{sc0-canonical-port:geom:tensor-sheaf}
\lean{EnvelopeProperty.CanonicalGeometry.bundleTensorIso}
\leanok
Actual sheaf isomorphisms \(\mathcal L\simeq\mathcal L'\)
and \(\mathcal N\simeq\mathcal N'\) induce an isomorphism
\(\mathcal L\otimes\mathcal N\simeq
\mathcal L'\otimes\mathcal N'\).
\end{lemma}
\begin{proof}
On each open, tensor the two linear isomorphisms. Tensoring their
inverses gives an inverse; naturality of both original maps makes
these isomorphisms commute with presheaf restrictions. Apply
sheafification to the two inverse presheaf maps. Their compositions
are still identities, giving the asserted sheaf isomorphism.
\end{proof}

\begin{lemma}[The actual tensor sheaf has the structure sheaf as unit]
\label{sc0-canonical-port:geom:tensor-sheaf-unit}
\uses{sc0-canonical-port:geom:tensor-sheaf}
\lean{EnvelopeProperty.CanonicalGeometry.bundleSheafificationIso,EnvelopeProperty.CanonicalGeometry.bundleTensorLeftUnitIso}
\leanok
Sheafifying the underlying presheaf of an actual sheaf recovers
that sheaf canonically. For every actual module sheaf
\(\mathcal L\), there is an isomorphism
\(\mathcal O_X\otimes\mathcal L\simeq\mathcal L\).
\end{lemma}
\begin{proof}
Maps from a sheafification to a sheaf are exactly maps from its
original presheaf. If the original presheaf already satisfies the
sheaf condition, its identity extends uniquely to a map inverse
to the sheafification unit; checking each composite uses this
same uniqueness. Objectwise, the maps
\(a\otimes s\mapsto a s\) and \(s\mapsto1\otimes s\)
are inverse module maps and commute with restriction. Thus the
tensor presheaf \(\mathcal O_X\otimes\mathcal L\) is
isomorphic to the underlying presheaf of \(\mathcal L\).
Sheafifying and using the first assertion proves the unit identity.
\end{proof}

\begin{lemma}[A global frame gives an actual rank-one atlas]
\label{sc0-canonical-port:geom:global-frame-line-atlas}
\uses{sc0-canonical-port:geom:local-line-bundle-frames}
\lean{EnvelopeProperty.CanonicalGeometry.trivialIntrinsicLineBundle,EnvelopeProperty.CanonicalGeometry.intrinsicLineBundleOfFrame}
\leanok
The structure sheaf is an intrinsic line bundle. More generally,
an actual sheaf isomorphism \(\mathcal L\simeq\mathcal O_X\)
gives a rank-one atlas for \(\mathcal L\).
\end{lemma}
\begin{proof}
Use the one-open cover consisting of \(X\) itself. The identity
frames the structure sheaf. For the general assertion, restrict
the given global isomorphism to this open. This is the required
actual frame and covers every point.
\end{proof}

\begin{theorem}[An actual affine linear trivialization gives a line sheaf]
\label{sc0-canonical-port:geom:affine-tilde-line-bundle}
\uses{sc0-canonical-port:geom:global-frame-line-atlas}
\lean{EnvelopeProperty.CanonicalGeometry.affineSheafFrameOfTrivialization,EnvelopeProperty.CanonicalGeometry.affineLineBundleOfTrivialization}
\leanok
Let \(A\) be a commutative ring and \(M\) an actual
\(A\)-module with a linear isomorphism \(M\simeq A\).
The actual associated module sheaf \(\widetilde M\) on
\(\operatorname{Spec}A\) has a global sheaf frame to the
structure sheaf and hence an intrinsic line-bundle atlas.
\end{theorem}
\begin{proof}
On every principal open \(D(r)\), localization transports the
linear isomorphism to \(M_r\simeq A_r\). These localized maps
commute with further localization and are inverse. Thus they
give inverse sheaf maps \(\widetilde M\simeq\widetilde A\).
The latter sheaf is the actual structure sheaf, since its
principal-open sections and restriction maps are \(A_r\) and
the standard localization maps. This proves the global frame.
Apply Lemma~\ref{sc0-canonical-port:geom:global-frame-line-atlas}.
\end{proof}

\begin{corollary}[Actual affine canonical and dual line sheaves]
\label{sc0-canonical-port:geom:affine-canonical-line-sheaves}
\uses{sc0-canonical-port:geom:affine-tilde-line-bundle,sc0-canonical-port:geom:canonical-frame-from-differential-basis,sc0-canonical-port:geom:anticanonical-dual-frame}
\lean{EnvelopeProperty.CanonicalGeometry.affineCanonicalLineBundle,EnvelopeProperty.CanonicalGeometry.affineAnticanonicalLineBundle}
\leanok
An actual basis of \(\Omega_{A/K}\) with \(n\) vectors
produces intrinsic line bundles on \(\operatorname{Spec}A\)
whose sheaves are the associated sheaves of the actual top
exterior differential module and its actual linear dual.
\end{corollary}
\begin{proof}
Theorems~\ref{sc0-canonical-port:geom:canonical-frame-from-differential-basis}
and \ref{sc0-canonical-port:geom:anticanonical-dual-frame} give their actual linear
trivializations to \(A\). Applying
Theorem~\ref{sc0-canonical-port:geom:affine-tilde-line-bundle} to each gives the
displayed sheaves with genuine global frames and atlases.
\end{proof}

\begin{corollary}[The polynomial chart produces actual canonical line sheaves]
\label{sc0-canonical-port:geom:polynomial-canonical-line-sheaves}
\uses{sc0-canonical-port:geom:affine-canonical-line-sheaves}
\lean{EnvelopeProperty.CanonicalGeometry.polynomialCanonicalLineBundle,EnvelopeProperty.CanonicalGeometry.polynomialAnticanonicalLineBundle}
\leanok
For \(A=K[X_0,\ldots,X_{n-1}]\), the actual canonical
exterior module and its actual dual give intrinsic line bundles
on \(\operatorname{Spec}A\) without a differential-basis
hypothesis.
\end{corollary}
\begin{proof}
The actual polynomial differential basis is produced in
Theorem~\ref{sc0-canonical-port:geom:polynomial-differential-basis}. Use it in
Corollary~\ref{sc0-canonical-port:geom:affine-canonical-line-sheaves}. Its empty
basis case also supplies the construction when \(n=0\).
\end{proof}

\begin{theorem}[A global sheaf frame produces every tensor-power frame]
\label{sc0-canonical-port:geom:globally-framed-tensor-powers}
\uses{sc0-canonical-port:geom:tensor-sheaf-powers,sc0-canonical-port:geom:tensor-sheaf-isomorphisms,sc0-canonical-port:geom:tensor-sheaf-unit,sc0-canonical-port:geom:global-frame-line-atlas}
\lean{EnvelopeProperty.CanonicalGeometry.bundleTensorPowerTrivialization,EnvelopeProperty.CanonicalGeometry.bundleTensorPowerLineBundle}
\leanok
If \(\mathcal L\) has an actual global frame
\(\mathcal L\simeq\mathcal O_X\), then every actual
\(\mathcal L^{\otimes n}\) has a global frame to
\(\mathcal O_X\), and hence a rank-one atlas.
\end{theorem}
\begin{proof}
For \(n=0\), the power is the structure sheaf. Suppose the
power \(n\) has the constructed frame. Tensor it with the
given frame of \(\mathcal L\), using
Lemma~\ref{sc0-canonical-port:geom:tensor-sheaf-isomorphisms}; then apply the
structure-sheaf unit isomorphism of
Lemma~\ref{sc0-canonical-port:geom:tensor-sheaf-unit}. This frames the recursive
power \(n+1\). Induction produces every frame, and
Lemma~\ref{sc0-canonical-port:geom:global-frame-line-atlas} gives its atlas.
\end{proof}

\begin{lemma}[Actual frame coefficients are linear and natural]
\label{sc0-canonical-port:geom:local-frame-coefficients}
\uses{sc0-canonical-port:geom:local-line-bundle-frames}
\lean{EnvelopeProperty.CanonicalGeometry.LocalBundleFrame.coefficients,EnvelopeProperty.CanonicalGeometry.LocalBundleFrame.coefficients_restrict}
\leanok
A local sheaf frame \(e\) on \(U\) gives on every
\(V\subset U\) an actual linear isomorphism
\(e_V:\Gamma(V,\mathcal L)\simeq\mathcal O_X(V)\).
For \(W\subset V\),
\(e_W(s|_W)=e_V(s)|_W\).
\end{lemma}
\begin{proof}
Evaluate the actual sheaf isomorphism on \(V\). Its evaluated
inverse proves the linear isomorphism. Naturality of the sheaf
map on the inclusion \(W\subset V\) is exactly the asserted
coefficient identity.
\end{proof}

\begin{theorem}[Overlap coefficients determine actual sections]
\label{sc0-canonical-port:geom:atlas-section-separation}
\uses{sc0-canonical-port:geom:local-frame-coefficients,sc0-canonical-port:geom:local-line-bundle-frames}
\lean{EnvelopeProperty.CanonicalGeometry.LineBundleAtlas.chartCoefficients,EnvelopeProperty.CanonicalGeometry.LineBundleAtlas.section_ext}
\leanok
For a genuine atlas \((U_i,e_i)\) and a section \(s\) on
\(U\), define its chart coefficient to be
\((e_i)_{U\cap U_i}(s|_{U\cap U_i})\).
If two actual sections \(s,t\) on \(U\) have equal chart
coefficients for every \(i\), then \(s=t\).
\end{theorem}
\begin{proof}
Each evaluated frame is injective, so equal coefficients imply
equal restrictions on \(U\cap U_i\). These intersections
cover \(U\), because the atlas covers \(X\). The separation
part of the sheaf axiom makes sections equal when their
restrictions on an open cover are equal, proving \(s=t\).
\end{proof}

\begin{lemma}[Actual rank-one transitions are units with a cocycle]
\label{sc0-canonical-port:geom:rank-one-transition-units}
\lean{EnvelopeProperty.CanonicalGeometry.rankOneAutomorphismUnit,EnvelopeProperty.CanonicalGeometry.rankOneAutomorphism_apply,EnvelopeProperty.CanonicalGeometry.sectionFrameTransition,EnvelopeProperty.CanonicalGeometry.sectionFrameTransition_coefficients,EnvelopeProperty.CanonicalGeometry.sectionFrameTransition_self,EnvelopeProperty.CanonicalGeometry.sectionFrameTransition_cocycle}
\leanok
For a commutative ring \(A\), every \(A\)-linear
automorphism \(h:A\to A\) is multiplication by the actual
unit \(h(1)\). For two linear frames \(e,f:M\simeq A\),
let \(u_{ef}=(f\circ e^{-1})(1)\). Then
\[
 f(s)=u_{ef}e(s),\quad u_{ee}=1,\quad
 u_{eg}=u_{fg}u_{ef}.
\]
\end{lemma}
\begin{proof}
Linearity gives \(h(a)=a h(1)\). The inverse map gives
\(h(1)h^{-1}(1)=1\), so \(h(1)\) is a unit.
Apply this identity to \(f\circ e^{-1}\) to get the
coefficient formula. The identity automorphism gives \(u_{ee}=1\).
For three frames, substitute the two coefficient identities
successively and evaluate at \(s=e^{-1}(1)\). This gives
the stated cocycle equality of units.
\end{proof}

\begin{lemma}[Actual sheaf transition units restrict as regular functions]
\label{sc0-canonical-port:geom:sheaf-frame-transition-units}
\uses{sc0-canonical-port:geom:local-frame-coefficients,sc0-canonical-port:geom:rank-one-transition-units}
\lean{EnvelopeProperty.CanonicalGeometry.localBundleFrameTransition,EnvelopeProperty.CanonicalGeometry.localBundleFrameTransition_coefficients,EnvelopeProperty.CanonicalGeometry.localBundleFrameTransition_cocycle,EnvelopeProperty.CanonicalGeometry.localBundleFrameTransition_restrict}
\leanok
For frames \(e,f\) on opens containing \(W\), their actual
transition \(u_{ef,W}\in\mathcal O_X(W)^\times\) satisfies
\(f_W(s)=u_{ef,W}e_W(s)\) and the same cocycle identity.
For \(Z\subset W\), \(u_{ef,Z}=u_{ef,W}|_Z\).
\end{lemma}
\begin{proof}
Apply Lemma~\ref{sc0-canonical-port:geom:rank-one-transition-units} to the
evaluated linear frames on \(W\); this gives the unit and
cocycle identities. Let \(s=e_W^{-1}(1)\). Naturality of
the coefficient maps gives \(e_Z(s|_Z)=1\) and
\(f_Z(s|_Z)=f_W(s)|_Z\). The first coefficient identity
on \(Z\) identifies its left side with \(u_{ef,Z}\);
the one on \(W\) identifies the restricted right side with
\(u_{ef,W}|_Z\). This proves the restriction formula.
\end{proof}

\begin{theorem}[A centered valuation gives frame-independent actual section values]
\label{sc0-canonical-port:geom:sheaf-frame-valuation-invariance}
\uses{sc0-canonical-port:geom:sheaf-frame-transition-units,sc0-canonical-port:fil:centered-unit-zero}
\lean{EnvelopeProperty.CanonicalGeometry.localBundleFrame_valuation_independent}
\leanok
Let \(W\) lie in two actual frame opens. Let
\(\iota:\mathcal O_X(W)\to F\) be a ring map to a field
with an additive real valuation \(v\), and assume
\(v(\iota(a))\geq0\) for every actual regular function
\(a\) on \(W\). For each actual section \(s\),
\(v(\iota(f_W(s)))=v(\iota(e_W(s)))\).
\end{theorem}
\begin{proof}
The transition formula gives \(f_W(s)=u e_W(s)\), where
\(u\) is an actual regular unit. Its value is zero because
the nonnegative values of \(u\) and \(u^{-1}\) add to zero,
as in Lemma~\ref{sc0-canonical-port:fil:centered-unit-zero}. The valuation product
law now proves the equality, including an infinite section value.
\end{proof}

\begin{definition}[The actual local section filtration]
\label{sc0-canonical-port:geom:actual-local-section-filtration}
\uses{sc0-canonical-port:geom:local-frame-coefficients,sc0-canonical-port:alg:valuation-ideal}
\lean{EnvelopeProperty.CanonicalGeometry.localBundleSectionFiltration,EnvelopeProperty.CanonicalGeometry.mem_localBundleSectionFiltration}
\leanok
For a frame \(e\) and a centered valued-field map \(\iota\)
as above, set
\[
 \mathcal F^a\Gamma(W,\mathcal L)
 =\{s:a\leq v(\iota(e_W(s)))\}.
\]
This is the actual \(\mathcal O_X(W)\)-submodule obtained
by pulling the valuation ideal back through the actual
coefficient linear isomorphism.
\end{definition}

\begin{theorem}[The actual section filtration is decreasing and independent of frame]
\label{sc0-canonical-port:geom:local-section-filtration-invariance}
\uses{sc0-canonical-port:geom:actual-local-section-filtration,sc0-canonical-port:geom:sheaf-frame-valuation-invariance}
\lean{EnvelopeProperty.CanonicalGeometry.localBundleSectionFiltration_antitone,EnvelopeProperty.CanonicalGeometry.localBundleSectionFiltration_frame_independent}
\leanok
If \(a\leq b\), then \(\mathcal F^b\subset\mathcal F^a\).
On a common open with the stated centered map, the whole
actual section filtration is independent of the chosen frame.
\end{theorem}
\begin{proof}
If a section value is at least \(b\), it is at least \(a\),
proving the inclusion. Theorem~\ref{sc0-canonical-port:geom:sheaf-frame-valuation-invariance}
makes the two frame values equal for every section. Comparing
the identical threshold inequalities proves equality of the
submodules for every real threshold.
\end{proof}

\begin{lemma}[The unselected coordinates of a submersive presentation]
\label{sc0-canonical-port:geom:submersive-free-coordinates}
\lean{EnvelopeProperty.CanonicalGeometry.submersiveFreeIndexEquiv}
\leanok
Let \(P\) be a finite submersive presentation of a commutative
\(K\)-algebra \(A\), with coordinates indexed by \(I\),
equations indexed by \(J\), and an injective selection
\(s:J\to I\) of Jacobian coordinates. Then
\(I\setminus s(J)\) has cardinality
\(\dim P=|I|-|J|\), and admits a bijection with
\(\{0,\ldots,\dim P-1\}\).
\end{lemma}
\begin{proof}
Injectivity of \(s\) makes its image have cardinality \(|J|\).
The image and its complement partition \(I\), so subtraction
gives the stated cardinality. Enumerating the finite complement
gives a bijection with its cardinality's standard finite index set.
\end{proof}

\begin{theorem}[The Jacobian relations produce an actual differential basis]
\label{sc0-canonical-port:geom:submersive-differential-basis}
\uses{sc0-canonical-port:geom:submersive-free-coordinates}
\lean{EnvelopeProperty.CanonicalGeometry.submersiveDifferentialBasis,EnvelopeProperty.CanonicalGeometry.submersiveDifferentialBasis_apply,EnvelopeProperty.CanonicalGeometry.submersiveDifferentials_finite_free,EnvelopeProperty.CanonicalGeometry.submersiveDifferentials_finrank}
\leanok
For the presentation \(P\) above, assume that the determinant
of its evaluated Jacobian matrix in the selected coordinates is
a unit in \(A\). The actual differentials of the unselected
coordinate images, enumerated as in
Lemma~\ref{sc0-canonical-port:geom:submersive-free-coordinates}, form an
\(A\)-basis of \(\Omega_{A/K}\).
Thus this module is finite free; if \(A\ne0\), its rank is
\(\dim P\).
\end{theorem}
\begin{proof}
If the presentation equations are \(F_j\), the differential
module is generated by the coordinate differentials, subject to
the relations
\[
 \sum_{i\in I}\overline{\frac{\partial F_j}{\partial X_i}}
 \,d\overline X_i=0\quad(j\in J).
\]
Indeed, every relation in the defining ideal is a sum of multiples
of the \(F_j\), whose derivative modulo that ideal is the
corresponding sum of multiples of \(dF_j\).
Separate selected and unselected coordinates. The selected block
of this relation matrix is invertible, by the adjugate formula
and the unit determinant. Multiplying the relations by its inverse
expresses each selected differential uniquely as a linear
combination of the unselected differentials.
The resulting quotient of the free coordinate module is therefore
the free module on the unselected coordinates: the map that keeps
the unselected coordinates and substitutes these formulas for the
selected ones kills exactly the submodule of relations. Its
inverse inserts the unselected differentials in the quotient.
This proves the basis assertion. Its finite index set gives
finite generation and freeness, and, over a nonzero ring, its
cardinality gives the stated rank.
\end{proof}

\begin{corollary}[A presentation with an invertible actual Jacobian]
\label{sc0-canonical-port:geom:jacobian-differential-basis}
\uses{sc0-canonical-port:geom:submersive-differential-basis}
\lean{EnvelopeProperty.CanonicalGeometry.jacobianPresentation_dimension,EnvelopeProperty.CanonicalGeometry.jacobianDifferentialBasis,EnvelopeProperty.CanonicalGeometry.jacobianDifferentialBasis_is_coordinate}
\leanok
Let a presentation of \(A\) have \(n+r\) coordinates and
\(r\) equations, with an injective choice of \(r\) Jacobian
coordinates. If its actual evaluated Jacobian determinant is a
unit, then \(\Omega_{A/K}\) has a basis indexed by
\(\{0,\ldots,n-1\}\). Each basis vector is
\(d\overline X_j\) for an actual unselected coordinate \(j\).
\end{corollary}
\begin{proof}
The presentation dimension is \((n+r)-r=n\). The unit
Jacobian condition turns this data into a submersive presentation.
Apply Theorem~\ref{sc0-canonical-port:geom:submersive-differential-basis} and
reindex its basis by this equality of dimensions. Its explicit
unselected-coordinate formula gives the last assertion.
\end{proof}

\begin{corollary}[Canonical and dual frames produced by the Jacobian]
\label{sc0-canonical-port:geom:jacobian-canonical-frames}
\uses{sc0-canonical-port:geom:jacobian-differential-basis,sc0-canonical-port:geom:canonical-frame-from-differential-basis,sc0-canonical-port:geom:anticanonical-dual-frame}
\lean{EnvelopeProperty.CanonicalGeometry.jacobianCanonicalFrame,EnvelopeProperty.CanonicalGeometry.jacobianAnticanonicalFrame,EnvelopeProperty.CanonicalGeometry.jacobian_canonical_anticanonical}
\leanok
For the presentation of Corollary~\ref{sc0-canonical-port:geom:jacobian-differential-basis},
the actual top exterior module
\(\bigwedge_A^n\Omega_{A/K}\) and its actual \(A\)-linear
dual have linear trivializations to \(A\), given by the produced
differential basis and its dual wedge. Both modules are finite
free; if \(A\ne0\), both have rank one.
\end{corollary}
\begin{proof}
The produced differential basis has \(n\) vectors. Its top
wedge is a one-element basis of the top exterior module by
Theorem~\ref{sc0-canonical-port:geom:canonical-frame-from-differential-basis}.
Evaluation on that vector is the dual trivialization of
Theorem~\ref{sc0-canonical-port:geom:anticanonical-dual-frame}. Their one-element
bases prove finite freeness and, for a nonzero ring, rank one.
\end{proof}

\begin{lemma}[Unit transitions between actual Jacobian frames]
\label{sc0-canonical-port:geom:jacobian-transition-unit}
\uses{sc0-canonical-port:geom:jacobian-differential-basis,sc0-canonical-port:geom:canonical-transition-unit}
\lean{EnvelopeProperty.CanonicalGeometry.jacobian_canonical_transition_det_isUnit}
\leanok
If two presentations of the same \(A/K\) have invertible
Jacobians and relative dimension \(n\), the determinant of
the coordinate change between their produced differential bases
is a unit in \(A\).
\end{lemma}
\begin{proof}
Both produced lists are actual bases of the same differential
module. Their coordinate-change matrices are inverse: expanding
one basis in the other and then back gives the identity on each
basis vector. Taking determinants makes the two determinants
inverse units, as also expressed in
Lemma~\ref{sc0-canonical-port:geom:canonical-transition-unit}.
\end{proof}

\begin{theorem}[Standard smoothness produces a differential basis]
\label{sc0-canonical-port:geom:standard-smooth-differential-basis}
\uses{sc0-canonical-port:geom:submersive-differential-basis}
\lean{EnvelopeProperty.CanonicalGeometry.exists_standardSmoothDifferentialBasis,EnvelopeProperty.CanonicalGeometry.standardSmoothDifferentialBasis,EnvelopeProperty.CanonicalGeometry.standardSmoothDifferentials_finite_free,EnvelopeProperty.CanonicalGeometry.standardSmoothDifferentials_finrank}
\leanok
Let \(A\) be an actual standard-smooth \(K\)-algebra of
relative dimension \(n\). Then \(\Omega_{A/K}\) admits a
basis indexed by \(\{0,\ldots,n-1\}\), obtained from a
submersive presentation of dimension \(n\).
In particular the module is finite free; if \(A\ne0\),
its rank is \(n\).
\end{theorem}
\begin{proof}
Standard smoothness of relative dimension \(n\) supplies a
finite submersive presentation with dimension \(n\), including
the actual unit Jacobian condition. Apply
Theorem~\ref{sc0-canonical-port:geom:submersive-differential-basis} to that
presentation and reindex by its dimension equality. This gives
the basis rather than assuming one. The finite-free and rank
conclusions follow from that basis as in the same theorem.
\end{proof}

\begin{corollary}[Canonical frames of actual standard-smooth algebras]
\label{sc0-canonical-port:geom:standard-smooth-canonical-frames}
\uses{sc0-canonical-port:geom:standard-smooth-differential-basis,sc0-canonical-port:geom:canonical-frame-from-differential-basis,sc0-canonical-port:geom:anticanonical-dual-frame}
\lean{EnvelopeProperty.CanonicalGeometry.standardSmoothCanonicalFrame,EnvelopeProperty.CanonicalGeometry.standardSmoothAnticanonicalFrame,EnvelopeProperty.CanonicalGeometry.standardSmooth_canonical_anticanonical}
\leanok
For an actual standard-smooth algebra of relative dimension
\(n\), the actual top differential module and its actual
linear dual have linear trivializations to \(A\). They are
finite free and, if \(A\ne0\), of rank one.
\end{corollary}
\begin{proof}
Apply the top-wedge and dual constructions to the produced
basis of Theorem~\ref{sc0-canonical-port:geom:standard-smooth-differential-basis}.
The top wedge and its dual are one-element bases, so their
coordinate maps give the trivializations and the asserted ranks.
\end{proof}

\begin{theorem}[Actual affine chart frames for a smooth scheme morphism]
\label{sc0-canonical-port:geom:smooth-scheme-chart-frames}
\uses{sc0-canonical-port:geom:standard-smooth-differential-basis,sc0-canonical-port:geom:standard-smooth-canonical-frames}
\lean{EnvelopeProperty.CanonicalGeometry.smoothRelativeDimension_chart_frames,EnvelopeProperty.CanonicalGeometry.smoothRelativeDimension_chart_basis_and_ranks}
\leanok
Let \(f:X\to Y\) be a scheme morphism smooth of relative
dimension \(n\), and let \(x\in X\). There are affine open
sets \(U\subset Y\) and \(V\subset X\), with \(x\in V\)
and \(V\subset f^{-1}(U)\), such that, for the actual
homomorphism \(B=\mathcal O_Y(U)\to A=\mathcal O_X(V)\),
\(\Omega_{A/B}\) has a produced basis of \(n\) vectors and
rank \(n\). The actual modules
\(\bigwedge_A^n\Omega_{A/B}\) and
\(\operatorname{Hom}_A(\bigwedge_A^n\Omega_{A/B},A)\)
are finite free of rank one and have linear frames to \(A\).
\end{theorem}
\begin{proof}
The affine-local characterization of smoothness of relative
dimension gives such \(U,V\) with the actual restriction
homomorphism standard smooth of relative dimension \(n\).
Theorem~\ref{sc0-canonical-port:geom:standard-smooth-differential-basis} supplies
the basis. The coordinate ring \(A\) is nonzero: the point
\(x\in V\) corresponds under the affine chart to a prime ideal,
and the zero ring has no prime ideals. Thus the rank assertion
and the finite-free rank-one assertions of
Corollary~\ref{sc0-canonical-port:geom:standard-smooth-canonical-frames} apply.
The same corollary gives the actual linear frames. These are
claims about the displayed coordinate modules, before any
comparison or gluing of differential sheaves.
\end{proof}

\begin{definition}[Actual module-presheaf restriction to an open site]
\label{sc0-canonical-port:geom:module-presheaf-over-open}
\lean{EnvelopeProperty.CanonicalGeometry.bundlePresheafOver}
\leanok
For an actual module presheaf \(P\) on \(X\) and an
actual open \(U\), let \(P|_U\) on the site of opens
\(V\subset U\) be the actual module \(P(V)\) with its
original restriction maps and actual ring \(\mathcal O_X(V)\).
\end{definition}

\begin{lemma}[The restricted sheafification unit is locally bijective]
\label{sc0-canonical-port:geom:restricted-sheafification-unit}
\uses{sc0-canonical-port:geom:module-presheaf-over-open}
\lean{EnvelopeProperty.CanonicalGeometry.bundleSheafificationOverMap,EnvelopeProperty.CanonicalGeometry.bundleSheafificationOverMap_underlying,EnvelopeProperty.CanonicalGeometry.bundleSheafificationOverMap_locally_bijective}
\leanok
Restricting the actual unit \(P\to P^\#\) gives the
actual map \(P|_U\to P^\#|_U\). Its underlying additive
presheaf map is the literal restriction of that unit.
It is locally bijective: equality of images is locally
equality of original sections, and each target section is
locally an image of an original section.
\end{lemma}
\begin{proof}
The restricted map applies the original unit on each actual
open \(V\subset U\), so its underlying map has the
claimed description. The unit of sheafification is locally
injective and locally surjective by its local equivalence
construction. Every cover witnessing these properties on
\(V\) has all its opens contained in \(V\subset U\),
so it is also a cover on the open site. Hence the same
local witnesses prove both properties after restriction.
\end{proof}

\begin{lemma}[Open-site sheafification inverts the actual restricted unit]
\label{sc0-canonical-port:geom:open-site-unit-inversion}
\uses{sc0-canonical-port:geom:restricted-sheafification-unit}
\lean{EnvelopeProperty.CanonicalGeometry.bundleSheafificationOverMap_becomes_iso,EnvelopeProperty.CanonicalGeometry.bundleOverCounitIso}
\leanok
Sheafifying the restricted unit on the open site makes it
an isomorphism. Sheafifying the underlying presheaf of an
actual sheaf on that site canonically recovers that sheaf.
\end{lemma}
\begin{proof}
A locally bijective map induces a sheaf isomorphism:
locally surjective representatives give a locally defined
inverse on each target section; local injectivity makes
these representatives independent of choices and compatible
on overlaps. The sheaf axiom glues them to a unique inverse.
Apply this to the restricted unit of the preceding lemma.
If the original object is already a sheaf, its unit is
inverted by the sheafification counit: both composites
are identities by the unique extension property for maps
from the original presheaf. This gives the second assertion.
\end{proof}

\begin{theorem}[Actual module sheafification commutes with open restriction]
\label{sc0-canonical-port:geom:sheafification-open-restriction}
\uses{sc0-canonical-port:geom:open-site-unit-inversion}
\lean{EnvelopeProperty.CanonicalGeometry.bundleOverSheafificationIso}
\leanok
There is an actual sheaf isomorphism
\(P^\#|_U\simeq(P|_U)^\#\).
\end{theorem}
\begin{proof}
Sheafify the restricted unit. By
Lemma~\ref{sc0-canonical-port:geom:open-site-unit-inversion}, it gives an
isomorphism from \((P|_U)^\#\) to the sheafification
of the underlying presheaf of \(P^\#|_U\). The latter
object is already a sheaf, so the canonical counit is an
isomorphism to \(P^\#|_U\). Compose these two actual
isomorphisms and invert, giving the asserted direction.
\end{proof}

\begin{lemma}[Actual objectwise tensor restriction]
\label{sc0-canonical-port:geom:tensor-presheaf-open-restriction}
\uses{sc0-canonical-port:geom:tensor-sheaf,sc0-canonical-port:geom:module-presheaf-over-open}
\lean{EnvelopeProperty.CanonicalGeometry.bundleOverTensorPresheaf,EnvelopeProperty.CanonicalGeometry.bundleOverTensorPresheafIso}
\leanok
Restricting the actual tensor presheaf of \(\mathcal L\)
and \(\mathcal N\) to \(U\) gives the actual objectwise
tensor presheaf of their restricted sheaves. The comparison
is the identity on each actual tensor module over
\(\mathcal O_X(V)\) for \(V\subset U\).
\end{lemma}
\begin{proof}
Both objects have at \(V\) exactly
\(\mathcal L(V)\otimes_{\mathcal O_X(V)}\mathcal N(V)\).
Both restriction maps send \(s\otimes t\) to the pure
tensor of its two restrictions. Pure tensors generate, so
these maps agree everywhere. Objectwise identities thus
form the asserted actual presheaf isomorphism.
\end{proof}

\begin{theorem}[Actual local frames tensor on their common open]
\label{sc0-canonical-port:geom:local-tensor-frame}
\uses{sc0-canonical-port:geom:sheafification-open-restriction,sc0-canonical-port:geom:tensor-presheaf-open-restriction,sc0-canonical-port:geom:tensor-sheaf-unit,sc0-canonical-port:geom:local-line-bundle-frames}
\lean{EnvelopeProperty.CanonicalGeometry.LocalBundleFrame.tensor}
\leanok
Actual frames of \(\mathcal L\) and \(\mathcal N\)
on the same actual open \(U\) construct an actual frame
of \(\mathcal L\otimes\mathcal N\) on \(U\).
\end{theorem}
\begin{proof}
First commute the tensor sheaf's sheafification with open
restriction by Theorem~\ref{sc0-canonical-port:geom:sheafification-open-restriction}.
Its defining restricted tensor presheaf is the tensor of
the two restricted sheaves by the preceding lemma.
Tensor the two actual sheaf frames to identify this
presheaf with the tensor of the two structure-sheaf
presheaves. Multiplication gives its unit isomorphism.
Sheafification and its counit recover \(\mathcal O_U\),
as in Lemma~\ref{sc0-canonical-port:geom:tensor-sheaf-unit}. Composing these
actual isomorphisms gives the required local frame.
\end{proof}

\begin{theorem}[A genuine local atlas produces every actual tensor-power atlas]
\label{sc0-canonical-port:geom:local-atlas-tensor-powers}
\uses{sc0-canonical-port:geom:local-tensor-frame,sc0-canonical-port:geom:tensor-sheaf-powers,sc0-canonical-port:geom:local-line-bundle-frames}
\lean{EnvelopeProperty.CanonicalGeometry.LocalBundleFrame.tensorPower,EnvelopeProperty.CanonicalGeometry.LineBundleAtlas.tensorPower,EnvelopeProperty.CanonicalGeometry.IntrinsicLineBundle.tensorPower}
\leanok
A frame on \(U\) constructs a frame of every actual
\(\mathcal L^{\otimes n}\) on \(U\). Consequently any
genuine line-bundle atlas constructs an atlas of every
power on its original open cover, and each power is an
actual intrinsic line bundle.
\end{theorem}
\begin{proof}
At \(n=0\), the actual power is the structure sheaf, with
its identity frame on \(U\). Given the frame at \(n\),
tensor it with the original frame by
Theorem~\ref{sc0-canonical-port:geom:local-tensor-frame}; the target is the
recursive actual power \(n+1\). Induction produces every
local frame. Apply this construction to each chart of
the original actual atlas and retain its actual opens,
index set, and covering proof. The constructed frames
then give an atlas and the corresponding intrinsic line
bundle for each power, without a global frame premise.
\end{proof}

\begin{definition}[Valuation threshold ideals]
\label{sc0-canonical-port:alg:valuation-ideal}
\lean{EnvelopeProperty.CanonicalGeometry.valuationIdeal,EnvelopeProperty.CanonicalGeometry.mem_valuationIdeal}
\leanok
Let \(A\) be a commutative ring, \(F\) a field, and
\(f:A\to F\) a unital ring homomorphism. Let
\(w:F\to\R\cup\{+\infty\}\) be an additive valuation, and
assume \(w(f(r))\geq0\) for every \(r\in A\). For a real
threshold \(a\), define
\[
 \mathfrak a_a(w,f)=\{r\in A:a\leq w(f(r))\}.
\]
This subset is an ideal of \(A\).
\end{definition}
\begin{proof}
The zero element has value \(+\infty\), so belongs to the subset.
If \(r,s\) belong, the additive valuation inequality gives
\[
 w(f(r+s))\geq\min\{w(f(r)),w(f(s))\}\geq a.
\]
Also \(w(f(-r))=w(-1)+w(f(r))=w(f(r))\), since
\(w(-1)=0\): multiplicativity applied to \((-1)^2=1\) gives
\(2w(-1)=0\). Thus the subset is an additive subgroup.
For \(t\in A\) and \(r\) in the subset, nonnegativity yields
\[
 w(f(tr))=w(f(t))+w(f(r))\geq0+a=a.
\]
It is therefore closed under multiplication by arbitrary ring
elements, proving ideal closure.
\end{proof}

\begin{lemma}[Decreasing ideal thresholds]
\label{sc0-canonical-port:alg:valuation-ideal-antitone}
\uses{sc0-canonical-port:alg:valuation-ideal}
\lean{EnvelopeProperty.CanonicalGeometry.valuationIdeal_antitone}
\leanok
If \(a\leq b\), then
\(\mathfrak a_b(w,f)\subseteq\mathfrak a_a(w,f)\).
\end{lemma}
\begin{proof}
For an element of \(\mathfrak a_b(w,f)\),
\(w(f(r))\geq b\geq a\), which is exactly membership in
\(\mathfrak a_a(w,f)\).
\end{proof}

\begin{lemma}[Whole-ring and proper thresholds]
\label{sc0-canonical-port:alg:valuation-ideal-top}
\uses{sc0-canonical-port:alg:valuation-ideal}
\lean{EnvelopeProperty.CanonicalGeometry.valuationIdeal_eq_top_iff,EnvelopeProperty.CanonicalGeometry.valuationIdeal_ne_top}
\leanok
For every real \(a\),
\[
 \mathfrak a_a(w,f)=A\quad\Longleftrightarrow\quad a\leq0.
\]
In particular, every positive threshold gives a proper ideal.
\end{lemma}
\begin{proof}
If \(a\leq0\), nonnegativity gives \(a\leq w(f(r))\) for every
\(r\), proving equality with \(A\). Conversely, equality makes
\(1\) belong to the ideal, and its value is
\(w(f(1))=w(1)=0\), so \(a\leq0\). For \(a>0\), the
contrapositive proves properness.
\end{proof}

\begin{lemma}[Products add element thresholds]
\label{sc0-canonical-port:alg:valuation-ideal-product-elements}
\uses{sc0-canonical-port:alg:valuation-ideal,sc0-canonical-port:fil:valuation-product}
If \(r\in\mathfrak a_a(w,f)\) and
\(s\in\mathfrak a_b(w,f)\), then
\(rs\in\mathfrak a_{a+b}(w,f)\).
\end{lemma}
\begin{proof}
Multiplicativity of \(f\) and \(w\) gives
\(w(f(rs))=w(f(r))+w(f(s))\geq a+b\), as in
Lemma~\ref{sc0-canonical-port:fil:valuation-product}. This inequality is the claimed
membership.
\end{proof}

\begin{lemma}[Products of valuation ideals]
\label{sc0-canonical-port:alg:valuation-ideal-product}
\uses{sc0-canonical-port:alg:valuation-ideal,sc0-canonical-port:alg:valuation-ideal-product-elements}
\lean{EnvelopeProperty.CanonicalGeometry.valuationIdeal_mul_le}
\leanok
For all real \(a,b\),
\[
 \mathfrak a_a(w,f)\,\mathfrak a_b(w,f)
 \subseteq\mathfrak a_{a+b}(w,f).
\]
\end{lemma}
\begin{proof}
The ideal product consists of finite sums of products \(r_i s_i\)
with \(r_i\in\mathfrak a_a(w,f)\) and
\(s_i\in\mathfrak a_b(w,f)\). Each product belongs to
\(\mathfrak a_{a+b}(w,f)\) by
Lemma~\ref{sc0-canonical-port:alg:valuation-ideal-product-elements}. Since this is an
ideal, it contains every finite sum of those products, proving the
inclusion.
\end{proof}

\begin{definition}[The positive valuation ideal]
\label{sc0-canonical-port:alg:positive-valuation-ideal}
\lean{EnvelopeProperty.CanonicalGeometry.positiveValuationIdeal,EnvelopeProperty.CanonicalGeometry.mem_positiveValuationIdeal}
\leanok
Under the ring, map, valuation, and nonnegativity hypotheses of
Definition~\ref{sc0-canonical-port:alg:valuation-ideal}, set
\[
 \mathfrak p(w,f)=\{r\in A:0<w(f(r))\}.
\]
This subset is an ideal of \(A\).
\end{definition}
\begin{proof}
The value at zero is \(+\infty>0\). If two values are strictly
positive, their minimum is strictly positive, so the valuation
inequality gives additive closure. Negation preserves the value.
For \(t\in A\), \(w(f(t))\geq0\), so multiplying an element
of strictly positive value by \(t\) again gives strictly positive
value. These are the ideal axioms.
\end{proof}

\begin{lemma}[Primality of the positive valuation ideal]
\label{sc0-canonical-port:alg:positive-valuation-ideal-prime}
\uses{sc0-canonical-port:alg:positive-valuation-ideal}
\lean{EnvelopeProperty.CanonicalGeometry.positiveValuationIdeal_isPrime}
\leanok
The ideal \(\mathfrak p(w,f)\) is prime.
\end{lemma}
\begin{proof}
It is proper because \(w(f(1))=0\). Suppose \(rs\) belongs
while neither \(r\) nor \(s\) belongs. Nonnegativity and failure
of strict positivity then give \(w(f(r))=w(f(s))=0\).
Multiplicativity would imply \(w(f(rs))=0\), contradicting
membership of \(rs\). Thus one factor belongs, proving primality.
\end{proof}

\begin{lemma}[Valuation-ideal quotient kernel]
\label{sc0-canonical-port:alg:valuation-quotient-kernel}
\uses{sc0-canonical-port:alg:valuation-ideal}
\lean{EnvelopeProperty.CanonicalGeometry.valuationIdeal_quotient_eq_zero_iff}
\leanok
For the canonical quotient homomorphism
\(q_a:A\to A/\mathfrak a_a(w,f)\),
\[
 \ker q_a=\mathfrak a_a(w,f).
\]
Consequently \(q_a(r)=0\) if and only if
\(a\leq w(f(r))\).
\end{lemma}
\begin{proof}
The quotient identifies two elements exactly when their difference
belongs to the ideal. Taking the second element to be zero shows
that the class of \(r\) vanishes exactly when
\(r\in\mathfrak a_a(w,f)\). This proves the kernel equality,
and the defining membership condition gives the displayed
valuation criterion. No finite-dimensionality of this quotient is
used in the argument.
\end{proof}

\begin{definition}[The finite-piece map into the valuation quotient]
\label{sc0-canonical-port:alg:finite-piece-quotient-map}
\uses{sc0-canonical-port:alg:valuation-ideal}
\lean{EnvelopeProperty.CanonicalGeometry.valuationPieceQuotientMap}
\leanok
Let \(K\) be a field, \(A\) a commutative \(K\)-algebra,
\(F\) a field extension of \(K\), and \(f:A\to F\) a unital
\(K\)-algebra homomorphism. Suppose that \(w(f(r))\geq0\)
for all \(r\in A\). Let \(V\) be a \(K\)-vector space with
linear \(\iota:V\to A\). Define
\[
 \pi_a:V\longrightarrow A/\mathfrak a_a(w,f),
 \qquad x\longmapsto[\iota(x)].
\]
This is a \(K\)-linear map, with no injectivity assumption on
\(f\) or \(\iota\).
\end{definition}

\begin{lemma}[The kernel of the actual finite-piece quotient map]
\label{sc0-canonical-port:alg:finite-piece-quotient-kernel}
\uses{sc0-canonical-port:alg:finite-piece-quotient-map,sc0-canonical-port:alg:valuation-quotient-kernel,sc0-canonical-port:fil:valuation-filtration}
\lean{EnvelopeProperty.CanonicalGeometry.valuationPieceQuotientMap_ker}
\leanok
In the setting of Definition~\ref{sc0-canonical-port:alg:finite-piece-quotient-map},
suppose in addition that \(w\) is trivial on \(K^*\). Then
\[
 \ker\pi_a=\mathcal F^a_w(V,f\circ\iota).
\]
\end{lemma}
\begin{proof}
The class \([\iota(x)]\) vanishes precisely when
\(\iota(x)\in\mathfrak a_a(w,f)\), by
Lemma~\ref{sc0-canonical-port:alg:valuation-quotient-kernel}. This is precisely the
inequality \(a\leq w(f(\iota(x)))\) defining the threshold
subspace. Equality follows for every vector \(x\).
\end{proof}

\begin{lemma}[Finite-piece rank and nullity]
\label{sc0-canonical-port:alg:finite-piece-rank-nullity}
\uses{sc0-canonical-port:alg:finite-piece-quotient-map,sc0-canonical-port:alg:finite-piece-quotient-kernel}
\lean{EnvelopeProperty.CanonicalGeometry.valuationPiece_colength_add_filtration}
\leanok
Under the hypotheses of Lemma~\ref{sc0-canonical-port:alg:finite-piece-quotient-kernel},
if \(V\) has finite dimension, then
\[
 \dim_K\operatorname{im}\pi_a
 +\dim_K\mathcal F^a_w(V,f\circ\iota)=\dim_K V.
\]
Only the image of this finite-dimensional piece is counted here;
the entire quotient \(A/\mathfrak a_a(w,f)\) need not be finite
 dimensional.
\end{lemma}
\begin{proof}
The first isomorphism theorem identifies
\(\operatorname{im}\pi_a\) with \(V/\ker\pi_a\).
Extend a basis of \(\ker\pi_a\) to a basis of \(V\).
The classes of the added basis vectors form a basis of this
quotient: they span by the extended basis decomposition, and a
linear combination vanishing in the quotient belongs to the kernel,
forcing all its added coordinates to vanish. Thus their number is
\(\dim_K V-\dim_K\ker\pi_a\). Substitute the kernel equality of
Lemma~\ref{sc0-canonical-port:alg:finite-piece-quotient-kernel} to obtain the formula.
\end{proof}

\begin{theorem}[Actual affine polynomial differential basis]
\label{sc0-canonical-port:geom:polynomial-differential-basis}
\lean{KaehlerDifferential.mvPolynomialBasis,KaehlerDifferential.mvPolynomialBasis_apply}
\mathlibok
For \(A=K[X_0,\ldots,X_{n-1}]\), the actual module
\(\Omega_{A/K}\) has \(A\)-basis
\((\mathrm dX_i)_{i=0}^{n-1}\).
\end{theorem}
\begin{proof}
Define the derivation \(D:A\to A^n\) by the tuple of formal
partial derivatives. The universal property of K\"ahler differentials
gives a linear map \(\Phi:\Omega_{A/K}\to A^n\) sending
\(\mathrm dP\) to \((\partial_iP)_i\). Define the opposite
linear map \(\Psi:(a_i)_i\mapsto\sum_i a_i\mathrm dX_i\).
Since \(\partial_iX_j\) is the Kronecker delta,
\(\Phi\Psi\) is the identity. The identity
\(\mathrm dP=\sum_i(\partial_iP)\mathrm dX_i\) holds for
constants and the generators \(X_j\), and is preserved under
addition and multiplication by the derivation product law.
Polynomial induction therefore proves it for every \(P\).
The elements \(\mathrm dP\) generate the differential module,
so \(\Psi\Phi\) is also the identity. This constructs the claimed
basis by an actual module isomorphism with \(A^n\).
\end{proof}
